% =============================================================================
%  CLAS12 Technical Note --- preamble
%
%  Document class, packages, and house-style definitions. This file is inputted
%  by template.tex before \begin{document}. Editing the note's style (fonts,
%  sizes, spacing, header/footer, front-matter macros) is done here.
%
%  The style mirrors the Pages template
%  templates/CLAS12_technical_note_template.pages:
%    * body        Times New Roman, 12 pt, ragged right, blank-line paragraphs
%    * headings    Helvetica bold, 14 pt (section) / 12.5 pt (subsection)
%    * title       Helvetica bold, 19 pt, centred
%    * author      Helvetica bold, 14 pt, centred
%    * affiliation Times New Roman italic, 12 pt, centred
%    * captions    Times New Roman italic, 10.5 pt, grey
%
%  Build (from this directory):  meson setup build && meson compile -C build
%  which runs latexmk + pdflatex.  See README.md.
% =============================================================================

\documentclass[12pt]{article}

% --- Page geometry: US Letter, 1 inch margins (matches the Pages template) ---
\usepackage[letterpaper,margin=1in]{geometry}

% --- Encoding -----------------------------------------------------------------
\usepackage[T1]{fontenc}
\usepackage[utf8]{inputenc}

% --- Fonts --------------------------------------------------------------------
%   mathptmx : Times-like serif for body text and maths (~ TimesNewRomanPSMT)
%   helvet   : Helvetica-like sans for the title block and section headings
%   courier  : Courier-like typewriter font for verbatim / code
\usepackage{mathptmx}
\usepackage[scaled=0.90]{helvet}
\usepackage{courier}

% --- Colours ------------------------------------------------------------------
%   captiongrey matches the Pages caption colour {26000,26000,26000}/65535 ~ 0.40
\usepackage{xcolor}
\definecolor{captiongrey}{gray}{0.40}

% --- Graphics, maths, tables, code -------------------------------------------
\usepackage{graphicx}
\usepackage{amsmath,amssymb}
\usepackage{booktabs}
\usepackage{tabularx,array}
\usepackage{placeins}
\usepackage{listings}
\lstset{basicstyle=\ttfamily\small,columns=fullflexible,breaklines=true,
        frame=single,rulecolor=\color{captiongrey}}

% --- Section headings: Helvetica bold, left aligned, house-style sizes --------
%   section    14   pt   (e.g. "1. Experiment Configuration")
%   subsection 12.5 pt   (e.g. "3.1 ...")
\usepackage{titlesec}
\titleformat{\section}
  {\sffamily\bfseries\fontsize{14}{17}\selectfont}{\thesection.}{0.5em}{}
\titleformat{\subsection}
  {\sffamily\bfseries\fontsize{12.5}{15}\selectfont}{\thesubsection}{0.5em}{}
\titleformat{\subsubsection}
  {\sffamily\bfseries\fontsize{12}{14}\selectfont}{\thesubsubsection}{0.5em}{}
\titlespacing*{\section}      {0pt}{16pt}{6pt}
\titlespacing*{\subsection}   {0pt}{12pt}{4pt}
\titlespacing*{\subsubsection}{0pt}{10pt}{4pt}

% --- Figure / table captions: Times italic, 10.5 pt, grey, "Figure 1. ..." ---
\usepackage{caption}
\DeclareCaptionFont{clascaption}%
  {\color{captiongrey}\itshape\fontsize{10.5}{12.5}\selectfont}
\captionsetup{font=clascaption,labelfont=clascaption,labelsep=period}

% --- Paragraph style: ragged right, blank-line separation, no indent ----------
\usepackage[document]{ragged2e}
\usepackage{parskip}

% --- Typographic tuning: better spacing, tolerate the odd long code token -----
\usepackage{microtype}
\setlength{\emergencystretch}{2em}

% --- Running header / footer --------------------------------------------------
\usepackage{fancyhdr}
\setlength{\headheight}{14.5pt}
\addtolength{\topmargin}{-2.5pt}

% --- Hyperlinks (load last) ---------------------------------------------------
\usepackage[hyperindex]{hyperref}
\hypersetup{
  colorlinks=true,
  linkcolor=black,
  citecolor=blue,
  urlcolor=blue,
  filecolor=magenta,
  pdfauthor={Maurizio Ungaro},
  pdftitle={CLAS12 CAD optimization}
}

% --- House-style front-matter helpers ----------------------------------------
\newcommand{\notetitle}[1]%
  {{\centering\sffamily\bfseries\fontsize{19}{23}\selectfont #1\par}}
\newcommand{\noteauthor}[1]%
  {{\centering\sffamily\bfseries\fontsize{14}{17}\selectfont #1\par}}
\newcommand{\noteaffiliation}[1]%
  {{\centering\itshape\fontsize{12}{15}\selectfont #1\par}}
\newcommand{\notenumberdate}[1]%
  {{\centering\normalfont\fontsize{11}{14}\selectfont #1\par}}
\newenvironment{noteabstract}%
  {\vspace{6pt}%
   {\centering\sffamily\bfseries\fontsize{14}{17}\selectfont Abstract\par}%
   \vspace{4pt}\begin{center}\begin{minipage}{0.86\textwidth}%
   \justifying\normalfont\normalsize}%
  {\end{minipage}\end{center}\vspace{6pt}}

\begin{document}
\pagestyle{fancy}
\fancyhf{}
\renewcommand{\headrulewidth}{0.4pt}
\renewcommand{\sectionmark}[1]{\markright{\thesection.\ #1}}
\fancyhead[L]{\footnotesize\slshape CAD optimization}
\fancyhead[R]{\footnotesize\slshape\rightmark}
\fancyfoot[L]{\small\slshape Jefferson Lab}
\fancyfoot[R]{\small\slshape M. Ungaro}
\fancyfoot[C]{\small\thepage}
\thispagestyle{plain}
\notetitle{Performance and geometry validation of optimized CLAS12 CAD models}
\vspace{18pt}
\noteauthor{Maurizio Ungaro}
\vspace{8pt}
\noteaffiliation{Jefferson Lab, 12,000 Jefferson Avenue, 23606 Newport News, VA, USA}
\vspace{6pt}
\notenumberdate{October 20, 2026}
\begin{noteabstract}
This study compares LTCC CAD optimization and sixteen separate GEMC3 beamline mesh studies with original
engineering STL files. The complete optimized LTCC model contains 93.9\% fewer triangles and bytes, with
median initialization CPU time reduced by 63.5\% and startup peak memory by 56.8\%. Nonoptical LTCC electron
transport is 1.29 times faster by CPU time, while hadron and optical workloads require 1.35--1.53 and 1.73
times as much CPU time; material and hierarchy changes and more complete transport after defect repairs may
contribute. The apex, downstream, and adapter/exit-shield studies reduce triangles by 45.4\%, 3.75\%, and
0.434\%, respectively, with apex startup wall time and peak memory reduced by 15.7\% and 17.4\%. The FT
adjuster consolidation reduces triangles by 98.74\% while matching an estimated reference union volume, and
the FT inner shield and flange assembly is reduced by 36.14\% with its enclosed pocket retained exactly and
sampled surface differences below 0.05 mm; the nine FT inner-shield bolts have 89.38\% fewer triangles,
matched copper volumes, and explicit clearance from the shield; the FT outer flange retains its rims with
60.79\% fewer triangles and twelve bolt holes filled; the repaired outer mount clears the flange with
11.44\% fewer triangles. These independent workloads do not establish detector-response equivalence or
complete overlap freedom, and remaining geometry and navigation diagnostics are reported explicitly. The
downstream FT inner shield is rebuilt with closed through-holes and retained longitudinal grooves and fully
smooth flange faces. The Moller cone has 68.32\% fewer triangles, filled surface holes, and a widened
opening that clears the FT inner shield; repaired-reference tracking separates decimation from source
defects. Subsequent socket relief resolves the remaining downstream tube contacts while preserving
placements. The separate cone tip has 63.58\% fewer triangles with retained topology and sharp faces.
The already-coarse FT flange permits a 2.17\% triangle reduction with unchanged tracking performance.
The inner Moller-cone extension restores a missing port wall and reduces triangles by 30.77\%,
with a separate repaired-reference tracking control.
The extension lead fill and casing retain their interfaces with fewer triangles; repeated isolated tracking
results are reported separately.
The Moller-cone extension tungsten add-on restores a missing port wall and reduces triangles by 56.80\%,
with a separate repaired-reference tracking control.
Local FT/cone joint reliefs clear the observed interfaces; tracking is benchmarked separately.
The separate HTCC transfer reduces three cone meshes by 49.20\% in triangles; six initialization checks pass,
with remaining non-manifold edges and a 0.053 mm entry-window contact recorded.
\end{noteabstract}
\setcounter{tocdepth}{1}
\tableofcontents
\clearpage
\FloatBarrier

\section{Scope and geometry changes}

The question is how much the optimized LTCC CAD implementation changes loading cost, memory use,
particle-transport cost, and geometry diagnostics. This is a software and geometry study, not a background-rate
prediction or a measurement of LTCC photon yield. No experimental data were collected.

The original data are taken from main revision ea3625f. The port uses the LTCC CAD changes in meson revision
4c7c6af. The later PMT absorption change in revision 3472ede is excluded, so PMT optical properties remain the
same in the compared geometries. The unrelated Meson build-system changes are also excluded; the same installed
executable runs every case. The source revisions are recorded in Refs.~\cite{baseline,optimized}.

The primary comparison restores all ten files from main's cadTempRemoved directory: S1-BW, S1-LW, S1-RW,
S1-TRB, NFrame, FrontPlate, Mount, Nose, BottomPlate, and Epoxy. Source volumes are placed in sector 3 and the
disabled frame and six-part nose copies are enabled in the other five sectors. The benchmark preserves the
legacy placement formulas and material assignments from ltcc\_frame.pl: source frame volumes are aluminum,
copied frames are stainless steel, nose pieces are aluminum, and Epoxy is G4\_CR39. No STL facet is repaired or
removed in the original baseline. This gives 90 frame/nose physical placements, compared with 24 in the
consolidated implementation, excluding the Winston cones.

The previous production configuration had this hardware disabled. Comparing it directly with the complete meson
model changes the amount of simulated hardware and answers a different question. Those initial measurements are
retained as supplementary context; the headline optimization results in this note compare the complete original
and consolidated models.

The consolidated implementation uses four hardware source volumes: backwall, rightwall, leftwall, and nose.
Sector 3 contains the source volumes; enabled sectors reuse them through CopyOf placements inside their LTCC
mothers. The port restores walls and nose hardware that were disabled in the original geometry, changes the
frame hierarchy and some materials, and replaces the previous chain of mother-volume subtractions with a
fitted, planar G4Trap intersection. The unmodified engineering drawings remain archived in the geometry source
repository. The restored original baseline keeps the legacy mothers and hierarchy unchanged.

The consolidated nose uses the implementation's effective material at density 3.86 g/cm\textsuperscript{3},
with mass fractions 0.9464 stainless steel, 0.0408 carbon, 0.0043 hydrogen, and 0.0085 oxygen. This study does
not establish material-budget or detector-response equivalence to the original model.

Three cases separate part of the optimization: Original uses the old data; Cones only replaces only WC\_S,
WC\_M, and WC\_L while retaining the restored original frame, mothers, and materials; Full port uses the new
geometry and all its optimized meshes. Cones only is a diagnostic configuration, not an implementation proposed
for deployment. Frame consolidation and mother changes remain combined in the Full port case and cannot be
assigned independent transport speedups from this comparison.

Figures~\ref{fig:stl-hardware}--\ref{fig:stl-cones} compare renderings of the original and optimized STL files
directly. The original includes all restored CAD hardware. Paired panels use identical orthographic cameras,
common bounds, and the same scale; no mesh repair, smoothing, or additional decimation is applied. The hardware
view applies the sector-3 placements, including the sector mother transform for the optimized model. The
resulting hardware transforms are checked against the restored original placements. Nose and cone close-ups use
their stored STL coordinate systems. Colors identify functional hardware groups, not materials. The images
expose simplification and changes in shape; they do not establish volume, material-budget, or detector-response
equivalence.

\begin{figure}[htbp]
\centering
\includegraphics[width=\textwidth]{figures/hardware_stl_comparison.pdf}
\caption{Sector-3 STL frame and nose hardware: complete original on the left and optimized on the right. Upper
panels show shaded front oblique views; lower panels show rear oblique views with triangle edges. The same
camera and scale are used within each row. The 15 original meshes contain 124,652 triangles; the four optimized
meshes contain 3,646. Winston cones, native LTCC mirrors, PMTs, mothers, and other sector copies are omitted.}
\label{fig:stl-hardware}
\end{figure}

\begin{figure}[p]
\centering
\includegraphics[width=\textwidth]{figures/nose_stl_comparison.pdf}
\caption{Nose consolidation in the stored engineering coordinates, with identical camera and scale: six
original meshes and 12,012 triangles on the left; one optimized mesh and 1,986 triangles on the right. Triangle
edges are shown; simplified details and differences in the rendered envelope are visible.}
\label{fig:stl-nose}
\end{figure}

\begin{figure}[p]
\centering
\includegraphics[width=\textwidth]{figures/cones_stl_comparison.pdf}
\caption{Winston-cone meshes with triangle edges: WC\_S, WC\_M, and WC\_L from top to bottom; original on the
left and optimized on the right. The camera, bounds, and scale match within each row. Original triangle counts
are 4,312, 3,076, and 2,578; each optimized mesh contains 1,500 triangles. These are individual source meshes,
not all physical cone placements.}
\label{fig:stl-cones}
\end{figure}

\FloatBarrier

\section{Simulation and measurement configuration}

All runs use the local gemc/dev module, which reports GEMC 6.2, and Geant4 10.7.4 on macOS arm64. A local CCDB
SQLite database prevents remote calibration access. Executable and STL SHA-256 hashes, the complete gcard,
commands, and platform information are retained with the results. References~\cite{gemc,cad} describe GEMC and
CAD import.

The physics list is STD + FTFP\_BERT, with Optical enabled only for the optical electron workload. Only LTCC is
present. RUNNO is explicitly 11 with the default geometry variation, enabling six LTCC sectors. The primary
momentum is 5 GeV/c, the vertex is at the origin, the polar angle spans 15--35 degrees, and the azimuth spans
the full circle. Both magnets are off; polarity is therefore not applicable. There is no target, no imposed
beam current, no luminosity normalization, no bunch-time window, and no background pileup. Each event contains
one primary particle. FTOn/FTOff and shielding configurations are not applicable to this isolated detector
test.

Initialization runs use N=0. Optical electron runs use 1,000 events; electron, proton, positive-pion, and
neutron runs without optical processes use 10,000 events each. Each timing workload has five repetitions with
seeds 12345--12349, rotating the order of the three geometry cases. The nonoptical tests use matched LUND
files: every geometry receives identical primary tracks for a given seed. The corrected optical study also uses
matched LUND primaries. Random histories of secondary particles can differ because the geometries and ensuing
transport differ.

Event output and the GUI are disabled. The benchmark applies identical default production cuts; no obsolete
frame-specific cut is used. No readout threshold or detector rate is evaluated. GEMC reports process CPU time
and its events-only CPU time separately. Its events-only timer excludes the first warmup event and scales the
remaining time to the full event count. The macOS time utility records wall time and maximum resident-set size.
Medians summarize five runs; observed ranges are provided where useful and are not confidence intervals. The
machine was not dedicated to the study, so wall time includes scheduling and system-load variation.

A separate CAD-only loading test removes the TEXT detector hierarchy and loads the active STL solids directly
into the world with standard aluminum, no sensitivities, no field, and N=0. The same executable and physics
selection are used. This artificial configuration measures loading of the mesh bundle without native LTCC
mothers, copies, or custom-material definitions; its coincident world placements are not a physical simulation
and are not overlap-checked.

Construction diagnostics use CHECK\_OVERLAPS=1, N=0, and seed 12345. GEMC performs placement construction
checks and invokes /geometry/test/run without overriding its sampling resolution. One diagnostic run per case
counts each G4Exception by code. Geant4 caps overlap reports per volume, and some warnings concern ineffective
sampling; the number of reports is not the number of all intersecting pairs. Runtime navigation warnings are
counted separately during event transport. No warning count is interpreted as proof of overlap-free geometry.

\FloatBarrier

\section{Mesh size and loading performance}

Table~\ref{tab:1} compares the active mesh inventory. Counts refer to unique STL source meshes selected by the
CAD gxml, not to repeated physical placements. All ten previously disabled meshes are active in the original
comparator. Nine frame meshes with 112,640 triangles become three meshes with 1,660 triangles (98.5\% fewer).
Six nose meshes with 12,012 triangles become one mesh with 1,986 triangles (83.5\% fewer). The three Winston
cones shrink from 9,966 to 4,500 triangles (54.8\% fewer), retaining three source meshes.

\begin{table}[htbp]

\caption{Active CAD inventory for default LTCC geometry.}

\label{tab:1}

\centering

\small

\setlength{\tabcolsep}{4pt}

\renewcommand{\arraystretch}{1.15}

\begin{tabularx}{\textwidth}{@{}Xrrr@{}}

\toprule

Quantity & Original & Cones only & Full port \\

\midrule

STL source meshes & 18 & 18 & 7 \\

Frame triangles & 112,640 & 112,640 & 1,660 \\

Nose triangles & 12,012 & 12,012 & 1,986 \\

Winston-cone triangles & 9,966 & 4,500 & 4,500 \\

All triangles & 134,618 & 129,152 & 8,146 \\

STL bytes & 6,732,412 & 6,459,112 & 407,888 \\

Frame/nose placements & 90 & 90 & 24 \\

TEXT geometry rows & 2,438 & 2,438 & 2,377 \\

Component rows & 668 & 668 & 662 \\

\bottomrule

\end{tabularx}

\end{table}

Table~\ref{tab:2} separates full-detector initialization from the artificial CAD-only loading test. LTCC
initialization CPU time decreases from 1.874 to 0.684 s (63.5\%); startup memory decreases from 176.75 to 76.31
MiB (56.8\%). CAD-only loading CPU time decreases by 92.6\%, demonstrating a loading benefit of the simplified
mesh bundle independently of the native mother redesign.

\begin{table}[htbp]

\caption{Median initialization measurements over five N=0 runs; 1 MiB = 1,048,576 bytes.}

\label{tab:2}

\centering

\small

\setlength{\tabcolsep}{4pt}

\renewcommand{\arraystretch}{1.15}

\begin{tabularx}{\textwidth}{@{}Xrrr@{}}

\toprule

Measurement & Original & Cones only & Full port \\

\midrule

LTCC CPU time (s) & 1.874 & 1.816 & 0.684 \\

LTCC wall time (s) & 2.592 & 2.530 & 1.388 \\

LTCC peak RSS (MiB) & 176.75 & 175.53 & 76.31 \\

CAD-only CPU time (s) & 1.321 & 1.271 & 0.098 \\

CAD-only wall time (s) & 2.066 & 2.027 & 0.835 \\

CAD-only peak RSS (MiB) & 165.97 & 162.89 & 66.22 \\

\bottomrule

\end{tabularx}

\end{table}

\FloatBarrier

\section{Electron and hadronic transport}

Table~\ref{tab:3} gives median events-only CPU time per 1,000 primaries. The nonoptical entries divide the
measured 10,000-event times by ten. Throughput factors are Original time divided by comparison time: values
above one indicate faster transport, and values below one indicate slower transport. Figure~\ref{fig:transport}
presents these ratios and the initialization ratio for direct comparison.

\begin{table}[htbp]

\caption{CPU seconds per 1,000 events and transport throughput factors; zero magnetic field.}

\label{tab:3}

\centering

\small

\setlength{\tabcolsep}{4pt}

\renewcommand{\arraystretch}{1.15}

\begin{tabularx}{\textwidth}{@{}Xrrrrr@{}}

\toprule

Workload & Original & Cones only & Full port & Cone factor & Port factor \\

\midrule

Electron + optical & 26.067 & 26.852 & 45.126 & 0.97 & 0.58 \\

Electron & 4.739 & 4.546 & 3.679 & 1.04 & 1.29 \\

Proton & 1.009 & 0.958 & 1.490 & 1.05 & 0.68 \\

Positive pion & 1.084 & 1.055 & 1.467 & 1.03 & 0.74 \\

Neutron & 0.575 & 0.601 & 0.877 & 0.96 & 0.66 \\

\bottomrule

\end{tabularx}

\end{table}

Optical electron transport: the cone-only substitution has a throughput factor of 0.97; the full port has a
factor of 0.58. Original per-1,000-event CPU times range from 25.549 to 29.228 s, and full-port times range
from 39.580 to 56.916 s.

Nonoptical electron transport: the cone-only substitution has a throughput factor of 1.04; the full port has a
factor of 1.29. Original per-1,000-event CPU times range from 4.207 to 6.261 s, and full-port times range from
3.476 to 6.536 s.

Proton transport: the cone-only substitution has a throughput factor of 1.05; the full port has a factor of
0.68. Original per-1,000-event CPU times range from 0.987 to 1.034 s, and full-port times range from 1.477 to
1.580 s.

Positive-pion transport: the cone-only substitution has a throughput factor of 1.03; the full port has a factor
of 0.74. Original per-1,000-event CPU times range from 1.060 to 1.170 s, and full-port times range from 1.444
to 1.546 s.

Neutron transport: the cone-only substitution has a throughput factor of 0.96; the full port has a factor of
0.66. Original per-1,000-event CPU times range from 0.543 to 0.651 s, and full-port times range from 0.842 to
0.908 s.

The cone-only comparison isolates the three substituted meshes. The full port additionally changes gas
containment, hardware coverage, placements, and materials. Optical paths and particle showers can therefore
differ, and faster loading does not imply faster transport or unchanged response. The measured changes include
both mesh-navigation cost and changes in physical transport; these data do not identify a unique cause. PMT
absorption was held fixed and cannot account for the difference between these cases.

Cone-only throughput ratios close to one represent small measured differences within the observed run
variation; these five repetitions do not establish a statistically significant speedup of the Winston-cone
substitution. The full-model timing also includes the revised materials and mothers, and does not separately
measure CPU cost per navigation query.

A separate matched 100-event electron diagnostic at seed 12345 enables tracking verbosity 1 and counts the
printed trace (Original: 18,271 track headers and 123,928 step rows; Cones only: 17,025 track headers and
115,895 step rows; Full port: 24,164 track headers and 203,087 step rows). Step rows include the initial row
for each track; they are not a direct count of navigation queries. These counts characterize the amount of
transport work in this sample. Verbose-run times are excluded from the timing tables, and this sample does not
isolate the contribution of each geometry change.

A separate matched 100-event proton diagnostic at seed 12345 enables tracking verbosity 1 and counts the
printed trace (Original: 6,135 track headers and 50,512 step rows; Cones only: 7,794 track headers and 62,163
step rows; Full port: 9,566 track headers and 99,568 step rows). Step rows include the initial row for each
track; they are not a direct count of navigation queries. These counts characterize the amount of transport
work in this sample. Verbose-run times are excluded from the timing tables, and this sample does not isolate
the contribution of each geometry change.

The increased CPU time is consistent with more complete transport in the optimized geometry. The original model
has invalid or rejected facets and repeated stuck tracks requiring corrective navigation pushes; these defects
may alter material crossings and suppress expected particle histories. In the matched diagnostic samples, the
optimized model prints 1.64 times as many electron step rows and 1.97 times as many proton step rows. Lower CPU
time in the defective original model may therefore reflect incorrect or incomplete transport rather than a
lower cost of equivalent physics. These observations support that interpretation but do not establish
causation: materials and mothers also change, and this study does not compare energy balance, particle
survival, or detector response event by event.

\begin{figure}[htbp]

\centering

\includegraphics[width=\textwidth]{figures/transport_comparison.pdf}

\caption{Original/comparison CPU-time ratios for initialization and particle transport. Bars use the ratio of
five-run medians. Transport uses the isolated default LTCC at RUNNO=11, 5 GeV/c primary momentum, and zero
magnetic field; only the optical electron workload enables optical processes. No readout threshold or
luminosity normalization applies. Ratios above one mean faster.}

\label{fig:transport}

\end{figure}

\FloatBarrier

\section{Warnings and geometry validation}

Table~\ref{tab:4} separates construction diagnostics from transport warnings. GeomVol1002 denotes placement
overlap or mother-containment reports; GeomSolids1001 records invalid-facet or surface-sampling warnings, and
GeomSolids1002 records rejected facets. GeomNav1002 denotes a stuck track for which the navigator attempts a
corrective push. GeomNav0003 denotes a stuck track abandoned with an event-abort request. Each exception header
is counted once rather than counting repeated wording within the warning.

\begin{table}[htbp]

\caption{Construction warning counts and total navigation warnings across five transport runs.}

\label{tab:4}

\centering

\small

\setlength{\tabcolsep}{4pt}

\renewcommand{\arraystretch}{1.15}

\begin{tabularx}{\textwidth}{@{}Xrrr@{}}

\toprule

Diagnostic & Original & Cones only & Full port \\

\midrule

GeomVol1002 (construction) & 1505 & 1515 & 1420 \\

GeomSolids1001 (construction) & 857 & 857 & 2 \\

GeomSolids1002 (construction) & 857 & 857 & 0 \\

All construction warning lines & 3219 & 3229 & 1422 \\

Mother containment reports & 792 & 792 & 517 \\

Sibling overlap reports & 713 & 723 & 903 \\

Optical electrons navigation & 405 & 395 & 0 \\

Electrons navigation & 3486 & 3562 & 0 \\

Protons navigation & 1946 & 1959 & 0 \\

Positive pions navigation & 1899 & 1923 & 0 \\

Neutrons navigation & 1810 & 1852 & 1 \\

GeomNav0003 (event-abort requests) & 0 & 2 & 0 \\

GeomVol1002 (transport) & 0 & 4 & 0 \\

HAD\_QGSPARTICIPANTS\_002 (transport) & 0 & 0 & 1 \\

\bottomrule

\end{tabularx}

\end{table}

The additional transport exception rows retain other diagnostics across all five workloads. The cone-only case
requests one event abort in optical electron seed 12347 and one in nonoptical electron seed 12349, both
following stuck tracks in the retained original frame geometry. Each reports that the track is abandoned after
25 zero steps. These runs exit normally and print timing footers, but that does not mean every event completed
its intended transport. Their timings are retained as measurements of this defective configuration, not
equivalent completed-event throughput. No GeomNav0003 event-abort request appears in the Original or Full port
cases. These mixed-workload totals are not event-normalized rates.

Original construction produces 3219 warning lines, while Full port produces 1422. Its diagnostic CPU time
changes from 237.66 to 257.80 s. Overlap-check timing is excluded from transport and initialization
comparisons. The stochastic construction check provides evidence of improvement where counts fall, but
remaining reports and unsampled intersections must be assessed separately.

Without overlap checking, each complete original import emits 857 invalid-facet warnings (GeomSolids1001) and
857 rejected-facet warnings (GeomSolids1002). The optimized import emits none of these warnings. Geant4
therefore excludes malformed facets from the original imported solids; the triangle counts in Table~\ref{tab:1}
describe the files rather than accepted Geant4 facets.

An independent binary-STL check finds exactly 857 zero-area triangles in the original files: 452 in S1-LW, 396
in S1-RW, six in Epoxy, and three in NFrame. The optimized active meshes contain no exact zero-area triangles.
This check uses the cross product of the stored single-precision vertex differences and tests for an exactly
zero vector; it does not establish that all meshes are watertight or free of near-degenerate facets.

Mother-containment reports change from 792 to 517; sibling-overlap reports change from 713 to 903. The changed
hierarchy alters what is tested, so these are counts of emitted reports rather than a one-to-one comparison of
all physical intersections. Native mirror and PMT/light-guide reports remain, and the new walls and nose also
appear in the diagnostic logs. Resolving the remaining geometry problems is separate from the performance port.

The ported geometry generator reproduces the packaged TEXT data exactly for all six LTCC variations. The shared
CI helpers select runs and variations; a missing mapping for run 4763 was corrected before comparison. For each
run, SQLite geometry and CAD are selected at the latest default row with run less than or equal to RUNNO. Row
counts are checked first, followed by the repository's row-by-row geometry comparison. All comparisons pass,
all CAD mother and mesh references resolve, and the generation preserves all non-LTCC SQLite rows. One-event
SQLite runtime smoke tests also complete for all six runs without G4Exception reports. Table~\ref{tab:5}
records the effective configurations and their row counts.

\begin{table}[htbp]

\caption{Ported model: validated run-dependent TEXT/SQLite geometry and CAD row counts.}

\label{tab:5}

\centering

\small

\setlength{\tabcolsep}{4pt}

\renewcommand{\arraystretch}{1.15}

\begin{tabularx}{\textwidth}{@{}rXrr@{}}

\toprule

RUNNO & TEXT variation & Geometry & CAD \\

\midrule

11 & default & 2377 & 7 \\

3029 & rga\_spring2018 & 1667 & 7 \\

4763 & rga\_fall2018 & 957 & 7 \\

6150 & rgb\_spring2019 & 957 & 7 \\

11323 & rgb\_spring2020 & 957 & 7 \\

15016 & rgm\_fall2021\_H & 1667 & 7 \\

\bottomrule

\end{tabularx}

\end{table}

\FloatBarrier
\section{Beamline port and validation: October 12 progress}

\subsection{Current port and validation status}

The current GEMC3 port includes all fifteen native beamline geometry and material variations: ten
standard run-mapped configurations and five detector-local configurations. It follows the active GEMC2
beamline dispatch; historical builders and disabled HD-ice daughters are not activated. All published
volumes use passive placement, and this passive system requires no sensitive digitization plugin or
optical surface.

Native geometry is published as beamline. The beamline\_cad companion contains the ten standard CAD
configurations plus DDVCS, while beamline\_transverse\_cad provides noWedges, forwardWedges, and
midpointWedges as separately selectable alternatives. Both native and CAD systems use the SQLite factory;
mesh coordinates remain in millimetres at scale one. The runtime inventory contains 51 STL files: 30
beamline meshes and seven meshes for each transverse alternative. The two apex meshes are optimized as
described below; the remaining 49 retain the original tessellation.

The semantic comparator passes all fifteen native geometry variations, fifteen material variations, and
fourteen CAD placement variations against the GEMC2 TEXT/GXML references. The ten standard runs also
match the named native and CAD entries in the GEMC2 SQLite reference. Seven identical duplicate CAD rows
at run 20000 are represented once per volume name; the DDVCS configuration is also available at its
detector-local run. This preserves the effective geometry without duplicating physical placements.

Saved construction smoke logs show completion for all fifteen native configurations and all three
transverse CAD alternatives with overlap checking enabled. These zero-event checks validate geometry
construction and mesh resolution, not particle transport or detector response. A separate default
beamline check after the apex optimization also exits normally with check\_overlaps=2, but records eight
GeomSolids1001 and nineteen GeomVol1002 exception blocks. These counts describe reported solid defects
and overlap/containment warnings, not unique defective meshes or overlapping pairs. The unchanged CAD
meshes still include surface-hole and facet-orientation diagnostics.

The current status is therefore a port with passing semantic comparisons and completed construction
checks, with residual geometry diagnostics retained for follow-up. Full-beamline transport performance,
overlap freedom, material-budget equivalence, magnetic-field operation, and full CLAS12 background
response equivalence have not been established by these checks. The isolated apex measurements below are
reported separately from this broader port validation.

\subsection{Apex shielding optimization and isolated benchmark}

The next optimization step covers two beamline apex shielding meshes. The lead fill is reduced from 2,880 to
1,440 triangles; the hub is renamed from Apex\_Shield\_PCAL\_Hub to Apex\_Shield\_Hub and reduced from
36,354 to 20,000 triangles. The two-file total falls from 39,234 to 21,440 triangles, a 45.4\% reduction.
Stored bytes fall from 1,961,868 to 1,072,168. All other beamline meshes remain unchanged. Materials,
coordinates, and passive placements are preserved in both repositories; all eleven CAD variations and eleven
run-dependent GEMC2 hub rows use the new name.

Topology-preserving quadric edge collapse keeps the lead opening and both hub components. Faces adjoining
hub edges with dihedral angles above five degrees are protected; no holes are filled and no component is
removed. Both originals and optimized meshes are closed and manifold, with no detected self-intersections or
exact zero-area triangles. Bidirectional checks include every vertex and approximately 200,000 edge/surface
samples in each direction. The largest sampled deviations are 0.0425 mm for the lead fill and 0.0209 mm for
the hub; these are sampled measurements, not proven global Hausdorff bounds. The lead volume changes by
-0.002287\%; the hub volume changes by -0.00000100\%. The original engineering meshes remain archived, and
the builder reproduces the optimized STL files byte for byte.

Table~\ref{tab:apex-meshes} summarizes the unchanged topology and reduced mesh inventory.

\begin{table}[htbp]
\centering
\caption{Apex engineering and optimized triangle counts; topology is preserved.}
\label{tab:apex-meshes}
\begin{tabular}{lrrrr}
\toprule
Mesh & Original & Optimized & Components & Genus \\
\midrule
Lead fill & 2,880 & 1,440 & 1 & 1 \\
Hub & 36,354 & 20,000 & 2 & 2 \\
\bottomrule
\end{tabular}
\end{table}

This follow-on benchmark uses one GEMC3 development executable dated October 11, 2026 and Geant4 11.4.3 on
macOS arm64. Its timings are independent of the GEMC2 LTCC measurements above and must not be compared
across engines. Only the apex pair is loaded, in air, with its original lead and carbon-steel materials.
Native beamline geometry and other CAD solids are omitted to isolate the substitution. Both cases use
FTFP\_BERT without optical processes, zero magnetic field, one thread, identical default production cuts,
and no GUI or event output. Fixed 5 GeV primaries originate at the target, uniform in solid angle from 0.25
to 3.5 degrees and over full azimuth. This aperture illuminates the forward shielding rather than reusing
the larger LTCC angular acceptance.

There are five alternating-order repetitions per case, with seeds 12345--12349 and matching LUND files for
each seed. Each electron, proton, positive-pion, and neutron run contains 200 events (1,000 events per
particle per geometry). This first beamline sample is smaller than the LTCC sample: showers in the thick
lead dominate its per-event cost. Tables report medians of the five runs. Initialization uses separate
zero-event runs with overlap checking disabled. Transport uses GEMC3's elapsed wall time since beamOn;
unlike the LTCC events-only CPU column, this is a wall-clock measure. Whole-process user plus system CPU
time and peak resident memory are also retained in the dataset. The original/optimized transport ratios are
ratios of medians, with no initialization subtraction.

Table~\ref{tab:apex-timing} gives loading, peak memory, and particle-transport measurements.

\begin{table}[htbp]
\centering
\caption{Median pair-only measurements over five repetitions. Transport entries are wall seconds
per 200 events; the final column is a ratio of medians.}
\label{tab:apex-timing}
\begin{tabular}{lrrr}
\toprule
Quantity & Original & Optimized & Original/optimized \\
\midrule
Initialization wall seconds & 1.151 & 0.970 & 1.19 \\
Loading peak RSS (MiB) & 186.09 & 153.64 & 1.21 \\
e- & 36.801 & 18.311 & 2.01 \\
proton & 9.329 & 5.174 & 1.80 \\
pi+ & 12.877 & 6.747 & 1.91 \\
neutron & 8.206 & 4.643 & 1.77 \\
\bottomrule
\end{tabular}
\end{table}

Median zero-event simulation wall time decreases by 15.7\% and loading peak resident memory decreases by
17.4\%. All 52 measured processes exit normally and print completion footers; all forty transport runs
report their 200 requested events. No Geant4 exception blocks are recorded in the pair-only loading,
transport, or separate check\_overlaps=2 construction samples. The semantic comparator passes all fifteen
native geometry/material variations and all fourteen CAD variations after the rename. Full beamline
construction checks are retained separately; remaining unoptimized meshes can still produce warnings. These
measurements establish the performance of this isolated sample, not full CLAS12 speed, overlap freedom, or
background-response equivalence. Stochastic shower histories may differ despite matched primaries and seeds.

\FloatBarrier

Figure~\ref{fig:apex-stl} compares the stored triangulations with matched cameras and bounds.

\begin{figure}[htbp]
\centering
\includegraphics[width=\textwidth]{figures/apex_stl_comparison.pdf}
\caption{Original (left) and optimized (right) apex lead-fill and hub STL surfaces. The views use
identical cameras and millimetre coordinates within each pair; edges show the stored triangulations.}
\label{fig:apex-stl}
\end{figure}

\FloatBarrier
\section{Downstream beam tube optimization}

The remaining downstream contacts recorded in this historical study are resolved by the subsequent socket
revision in Section~\ref{sec:downstream-joints}.

The next beamline step covers the two downstream tube segments. Their engineering tessellations are already
coarse, so reductions are modest: 5,146 to 5,000 and 672 to 600 triangles. Together they decrease from 5,818
to 5,600 triangles (3.75\%) and from 291,068 to 280,168 bytes. Names, stainless-steel materials, coordinate
frames, and placements are retained in GEMC3 and GEMC2; no CAD database row is renamed or moved.

Segment 1 originally contains seven disconnected surface components, 392 boundary edges, fourteen boundary
loops, and 1,021 faces flagged by the mesh self-intersection filter. Generic hole closing caps the bore and
is unsuitable. The repair registers the seven inner/outer boundary-ring pairs on matching angular grids
while preserving their measured radii and axial positions, then adds 392 annular faces. All 392 seam
vertices are protected during decimation. The output is a single closed, manifold tube of genus one; the
passage remains open. Segment 2 retains its original single-component, genus-one topology. Neither output
has detected self-intersections or exact zero-area triangles.

Table~\ref{tab:downstream-meshes} summarizes the mesh inventory.

\begin{table}[htbp]
\centering
\caption{Downstream triangle counts and surface-component changes.}
\label{tab:downstream-meshes}
\begin{tabular}{lrrrr}
\toprule
Mesh & Original & Optimized & Original components & Output components \\
\midrule
Segment 1 & 5,146 & 5,000 & 7 & 1 \\
Segment 2 & 672 & 600 & 1 & 1 \\
\bottomrule
\end{tabular}
\end{table}

The largest bidirectionally sampled surface differences are 0.7454 mm and 0.4574 mm, mainly around short
connector and flange features. Bounds change by at most 0.011 mm. The 0.05 mm apex sampling limit does not
apply to these tubes. Every original/output vertex and approximately 200,000 edge/surface samples per
direction are checked; these sampled maxima are not proven global error bounds. Segment 1 has no
well-defined original volume because its surface is open. Decimation changes its annularly repaired
reference volume by -0.01632\%; this does not establish original-to-repaired mass equivalence. Segment 2
changes by -0.00170\%. The builder reproduces both outputs byte for byte from archived originals.

The pair-only benchmark uses the recorded GEMC3 executable and Geant4 stack, FTFP\_BERT without optical
processes, air, zero magnetic field, one thread, and identical default production cuts. GUI and event output
are disabled. Matched 5 GeV LUND primaries start at the target, uniform in solid angle from 0.1 to 1.3
degrees and over full azimuth. The smaller forward aperture illuminates the downstream tubes. Five
alternating-order repetitions use seeds 12345--12349 and 200 events per electron, proton, positive-pion, or
neutron run. Separate zero-event runs measure startup without overlap checking; elapsed wall time since
beamOn measures transport. CPU and peak RSS are retained separately. Absolute timings are not pooled with
the apex or LTCC workloads.

Table~\ref{tab:downstream-timing} gives the separate particle-transport measurements.

\begin{table}[htbp]
\centering
\caption{Median transport wall seconds per 200 events over five repetitions.}
\label{tab:downstream-timing}
\begin{tabular}{lrrr}
\toprule
Particle & Original & Optimized & Original/optimized \\
\midrule
e- & 13.008 & 11.401 & 1.14 \\
proton & 2.109 & 1.901 & 1.11 \\
pi+ & 2.318 & 2.161 & 1.07 \\
neutron & 1.340 & 1.204 & 1.11 \\
\bottomrule
\end{tabular}
\end{table}

Median startup simulation wall time changes from 0.6924 to 0.6925 s; loading peak RSS changes from 123.27 to
123.06 MiB. These small initialization differences should not be interpreted as a substantial startup or
memory improvement. The original/optimized median transport wall-time ratios range from 1.07 to 1.14. All 52
processes exit normally, and all forty transport runs report their 200 requested events. No event-abort
exception is recorded.

Table~\ref{tab:downstream-warnings} quantifies the retained geometry and navigation diagnostics.

\begin{table}[htbp]
\centering
\caption{Downstream construction and transport warning counts.}
\label{tab:downstream-warnings}
\begin{tabular}{lrr}
\toprule
Diagnostic & Original & Optimized \\
\midrule
Surface-hole warnings, five loading runs & 5 & 0 \\
e- stuck tracks, five runs & 10 & 14 \\
proton stuck tracks, five runs & 0 & 0 \\
pi+ stuck tracks, five runs & 0 & 4 \\
neutron stuck tracks, five runs & 0 & 1 \\
Placement reports, overlap sample & 1 & 3 \\
\bottomrule
\end{tabular}
\end{table}

The repair eliminates the original surface-hole import warning. Placement overlap reports and stuck-track
diagnostics remain, as quantified below. A completion footer and requested event count do not prove that
every track followed the intended geometry. The original mesh lacks closed surfaces, and the repaired model
changes the represented material boundaries. Therefore transport timing differences cannot be assigned
solely to triangle-count reduction or used as proof of faster correct transport. The unchanged engineering
placement overlap and the thin, coarsely tessellated connector surfaces require further geometry/navigation
work. Detector response and full-experiment performance are outside this study.

\FloatBarrier
Figure~\ref{fig:downstream-stl} compares the stored triangulations.

\begin{figure}[htbp]
\centering
\includegraphics[width=\textwidth]{figures/downstream_stl_comparison.pdf}
\caption{Original (left) and optimized (right) downstream tube segments in millimetre coordinates.
Both models retain the through-bore; the paired views use the same camera and scale.}
\label{fig:downstream-stl}
\end{figure}

\FloatBarrier
\section{Adapter and torus exit shield}

The remaining downstream contacts recorded in this historical study are resolved by the subsequent socket
revision in Section~\ref{sec:downstream-joints}.

The third beamline study covers the adapter and torus exit-shield casing and lead fill. The reducer is
renamed Beamline\_Adapter\_1 in GEMC3 and GEMC2, including all eleven source CAD variations, generated
runtime GXML, and eleven run-dependent SQLite CAD rows. Every other database field, material, coordinate
frame, and placement is retained; custom GXML must use the new adapter name to match the STL filename.

These three meshes are already coarse, closed, manifold, and free of detected self-intersections and
zero-area triangles. The adapter decreases from 528 to 526 triangles and the casing from 1,392 to 1,384. The
lead fill remains byte-identical at 384 triangles, with both closed components retained: even the smallest
trial collapse displaced its surface by more than 1 mm near tight mating clearances. Its exact geometry was
explicitly retained after reviewing this constraint. The group decreases from 2,304 to 2,294 triangles
(0.434\%) and from 115,452 to 114,952 bytes.

Table~\ref{tab:exit-meshes} summarizes the mesh inventory.

\begin{table}[htbp]
\centering
\caption{Adapter and exit-shield triangle counts.}
\label{tab:exit-meshes}
\begin{tabular}{lrr}
\toprule
Mesh & Original & Optimized \\
\midrule
Adapter & 528 & 526 \\
Exit casing & 1,392 & 1,384 \\
Exit lead fill & 384 & 384 \\
\bottomrule
\end{tabular}
\end{table}

For the casing, 432 faces having at least one vertex within 1 mm of the lead fill are protected. The saved
output is checked to retain every protected triangle at exactly the original coordinates. All three meshes
retain their bounds, component counts, and genus; no hole filling or reconstruction is used. The adapter and
casing have maximum bidirectionally sampled surface differences of 0.4559 and 0.5002 mm, below the 0.51 mm
sampling guard for this group, and volume changes of +0.000368\% and +0.000188\%. The guard is specific to
this coarse group and does not establish the apex study's 0.05 mm criterion. Every vertex and approximately
200,000 edge/surface samples per direction are checked; the sampled maxima are not proven global bounds. The
builder reproduces all three outputs byte for byte from the archived originals.

The benchmark isolates these three volumes with their unchanged aluminum, stainless-steel, and lead
materials in air. The recorded GEMC3 executable and Geant4 stack run FTFP\_BERT without optical processes,
zero magnetic field, one thread, default production cuts, and no GUI or event output. Matched 5 GeV LUND
primaries start at the target, uniform in solid angle over 0.25--2.6 degrees and full azimuth. Five
alternating-order repetitions use seeds 12345--12349 and 200 events per electron, proton, positive-pion, or
neutron run. Separate zero-event runs measure startup without overlap checking; time since beamOn measures
transport. CPU time and peak RSS are retained separately. Absolute timings are not pooled with either
earlier beamline study or LTCC.

Table~\ref{tab:exit-timing} gives the separate transport measurements.

\begin{table}[htbp]
\centering
\caption{Median transport wall seconds per 200 events over five repetitions.}
\label{tab:exit-timing}
\begin{tabular}{lrrr}
\toprule
Particle & Original & Optimized & Original/optimized \\
\midrule
e- & 10.750 & 10.443 & 1.029 \\
proton & 3.451 & 3.398 & 1.016 \\
pi+ & 4.363 & 4.269 & 1.022 \\
neutron & 3.421 & 3.293 & 1.039 \\
\bottomrule
\end{tabular}
\end{table}

Median startup simulation wall time changes from 0.8844 to 0.8707 s, and peak loading RSS from 125.23 to
125.16 MiB. The original/optimized median transport wall-time ratios range from 1.016 to 1.039; individual
seed-paired ratios span 1.015--1.052. With only ten removed triangles and five repetitions, these timings do
not establish a substantial performance gain. All 52 processes exit normally, and all forty transport runs
report the 200 requested events.

Table~\ref{tab:exit-warnings} gives the sampled diagnostic counts.

\begin{table}[htbp]
\centering
\caption{Geant4 exception counts over all 52 isolated group processes.}
\label{tab:exit-warnings}
\begin{tabular}{lrr}
\toprule
Diagnostic & Original & Optimized \\
\midrule
GeomSolids1001 & 0 & 0 \\
GeomNav1002 & 0 & 0 \\
GeomNav0003 & 0 & 0 \\
GeomVol1002 & 0 & 0 \\
\bottomrule
\end{tabular}
\end{table}

No Geant4 exceptions are recorded in the isolated loading, transport, or overlap samples for either case.
This is a limited sample, not proof of complete overlap freedom or detector-response equivalence. Full
default beamline construction also completes with the renamed adapter and no missing mesh; its seven solid
warnings and twenty overlap reports are unchanged from the preceding downstream state. All fifteen native
geometry/material variations and fourteen CAD variations pass the semantic comparison against GEMC2.

\FloatBarrier
Figure~\ref{fig:exit-stl} compares the stored triangulations.

\begin{figure}[htbp]
\centering
\includegraphics[width=0.84\textwidth]{figures/exit_stl_comparison.pdf}
\caption{Original (left) and final (right) adapter and torus exit-shield STL surfaces in millimetres.
Each pair shares camera and bounds; the lead fill is retained exactly.}
\label{fig:exit-stl}
\end{figure}

\FloatBarrier
\section{FT adjuster material model}

The fourth beamline study replaces adjuster1, adjuster2, and adjuster3 with one ft\_adjuster mesh. Here the
intended optimization prioritizes the amount of stainless steel over detailed shape. All three source pieces
share the same material, color, mother, position, and passive rotation; the replacement retains these
physical attributes and merges the three placement rows into one in each of the nine applicable CAD
variations. GEMC2's corresponding 27 SQLite CAD rows become nine rows, with every other CAD row retained.

The originals contain 1,687, 424, and 1,689 triangles. The first and third meshes each lack a nine-edge
surface patch, and the assembled pieces intersect. Their original material volume is therefore not
well-defined directly from the raw shells. Small sealing patches define a reference, followed by Boolean
union to avoid counting overlapping steel twice. These temporary reference meshes retain local
intersections; none is published as runtime geometry. The reference defines an estimated material budget,
not a certified engineering mass.

The sealed parts sum to 25203.967 mm$^3$, and their reference union contains 24922.570 mm$^3$. Pairwise
intersection volumes are 140.391 and 141.004 mm$^3$; the first and third pieces have disjoint longitudinal
bounds. Independent inclusion-exclusion reproduces the union volume within 0.00001\%. This subtracts about
1.12\% from the sum of overlapping parts. Preserving their unqualified sum would retain double-counted
material.

Table~\ref{tab:adjuster-material} gives the material-volume comparison.

\begin{table}[htbp]
\centering
\caption{Adjuster material volumes in cubic millimetres.}
\label{tab:adjuster-material}
\begin{tabular}{lr}
\toprule
Reference or model & Volume \\
\midrule
Sealed parts sum & 25203.966821 \\
Inclusion-exclusion union & 24922.571202 \\
Reference Boolean union & 24922.569617 \\
Final saved STL & 24922.562549 \\
\bottomrule
\end{tabular}
\end{table}

The final model is a closed 12-sided prism with radius 7.383895 mm and length 152.370361 mm, centered on the
reference union's centroid. It retains the longitudinal span and approximate off-axis material location. Its
saved float32 STL volume is 24922.562549 mm$^3$, differing from the reference by -0.00002836\%. The
unchanged stainless-steel material retains its density, so this relative volume match also matches the
reference mass. The output has one closed manifold component of genus zero, with no detected
self-intersections or zero-area faces, and reproduces byte for byte from archived originals. Shape detail,
holes, threads, and the original material-path distribution are deliberately changed; detector-response
equivalence is not established.

The runtime model decreases from three files and 3,800 triangles to one file and 48 triangles (98.74\%
fewer). STL bytes decrease from 190,252 to 2,484 (98.69\%). One placement replaces three. This is a
material-volume substitution rather than a surface-tolerance decimation.

Table~\ref{tab:adjuster-meshes} summarizes the consolidation.

\begin{table}[htbp]
\centering
\caption{Original adjuster trio and simplified material model.}
\label{tab:adjuster-meshes}
\begin{tabular}{lrrr}
\toprule
Model & Files & Triangles & Bytes \\
\midrule
Original trio & 3 & 3,800 & 190,252 \\
ft adjuster & 1 & 48 & 2,484 \\
\bottomrule
\end{tabular}
\end{table}

The benchmark compares the original raw trio with the merged model in air using the recorded GEMC3
executable and Geant4 stack, FTFP\_BERT without optical processes, zero magnetic field, one thread, default
production cuts, and no GUI or event output. Fixed 5 GeV LUND primaries start at the target, uniform in
solid angle over 2.5--3.3 degrees and azimuth -95-- -85 degrees. The narrow azimuth targets the off-axis
adjuster instead of diluting the sample over full azimuth. Five alternating-order repetitions use seeds
12345--12349 and 200 events for each particle. Zero-event runs measure startup without overlap checking;
time since beamOn measures transport. Both fixtures share physical placement attributes, while the original
contains three solids and the final contains one. The benchmark accepts a generated database from either
geometry state.

Table~\ref{tab:adjuster-timing} gives median transport wall seconds per 200 events.

\begin{table}[htbp]
\centering
\caption{Median transport wall seconds over five repetitions.}
\label{tab:adjuster-timing}
\begin{tabular}{lrrr}
\toprule
Particle & Original & Final & Original/final \\
\midrule
e- & 4.871 & 3.124 & 1.559 \\
proton & 0.882 & 0.784 & 1.126 \\
pi+ & 0.903 & 0.805 & 1.121 \\
neutron & 0.699 & 0.597 & 1.170 \\
\bottomrule
\end{tabular}
\end{table}

Median startup simulation wall time changes from 0.6875 to 0.6803 s; peak loading RSS changes from 119.55 to
114.31 MiB. The median original/optimized transport wall-time ratios range from 1.121 to 1.559. CPU time,
process wall time, memory, and raw diagnostics are retained separately. All 52 processes exit normally and
all forty transport runs report the 200 requested events.

Table~\ref{tab:adjuster-warnings} gives Geant4 exception counts over all 52 processes.

\begin{table}[htbp]
\centering
\caption{Geant4 exception counts for the isolated adjuster samples.}
\label{tab:adjuster-warnings}
\begin{tabular}{lrr}
\toprule
Diagnostic & Original & Final \\
\midrule
GeomSolids1001 & 52 & 0 \\
GeomNav1002 & 40 & 0 \\
GeomNav0003 & 0 & 0 \\
GeomVol1002 & 8 & 0 \\
\bottomrule
\end{tabular}
\end{table}

The repeated diagnostic counts are retained explicitly. Full default beamline construction completes with no
missing mesh and 5 solid warnings and 10 overlap reports, compared with seven and twenty before this merge.
These counts describe the check sample, not proof of complete overlap freedom. All fifteen native
geometry/material variations and fourteen CAD variations pass the semantic comparison against GEMC2.
Material redistribution, closure repairs, removal of internal overlaps, and fewer solids all affect
transport; timing changes cannot be assigned solely to triangle count or taken as evidence of unchanged
detector response. A normal completion footer does not certify every track in the original imperfect
geometry.

\FloatBarrier
Figure~\ref{fig:adjuster-stl} shows the deliberate shape simplification.

\begin{figure}[htbp]
\centering
\includegraphics[width=\textwidth]{figures/adjuster_stl_comparison.pdf}
\caption{Engineering adjuster trio (left) and the volume-matched ft\_adjuster prism (right).
Matched cameras and bounds use original STL coordinates; the common CAD transform is retained.}
\label{fig:adjuster-stl}
\end{figure}

\FloatBarrier
\section{FT inner shield and flange}

The fifth beamline study renames innerShieldAndFlange to FTInnerShieldAndFlange and optimizes its
engineering tessellation. The same rename applies to ten source and generated runtime CAD variations and ten
GEMC2 SQLite CAD rows. The tungsten-alloy material, density, color, mother, position, passive rotation,
coordinate frame, and overall bounds remain unchanged. Custom CAD placements must use the new name to match
the STL filename.

The original has 34,448 triangles, two closed manifold surface components, aggregate genus 17, and no
detected self-intersections or exact zero-area faces. The small inward-oriented component is an enclosed
pocket with signed volume -9.663980 mm$^3$. All 474 of its triangles are protected and verified in the saved
STL at exactly the original coordinates. Only the larger component is decimated, using topology-preserving
quadric collapse. No holes are filled, components removed, or surfaces reconstructed.

The saved output contains 22,000 triangles (36.14\% fewer) and 1,100,084 bytes instead of 1,722,484. It
retains both components, aggregate genus, and bounds; no self-intersections or zero-area triangles are
detected. Bidirectional vertex, edge, and surface sampling measures a maximum difference of 0.0481 mm, below
the 0.05 mm guard. Every vertex and approximately 200,000 edge/surface samples per direction are checked;
these sampled maxima are not proven global error bounds. Total signed volume changes by -0.00004391\%. A
separate component integration gives -0.00004114\% for the main component and no change for the pocket. The
builder reproduces the optimized STL byte for byte from the archived original.

Table~\ref{tab:inner-shield-meshes} summarizes the mesh inventory.

\begin{table}[htbp]
\centering
\caption{Original and optimized FT inner shield and flange.}
\label{tab:inner-shield-meshes}
\begin{tabular}{lrr}
\toprule
Quantity & Original & Optimized \\
\midrule
Triangles & 34,448 & 22,000 \\
Bytes & 1,722,484 & 1,100,084 \\
Closed surface components & 2 & 2 \\
Aggregate genus & 17 & 17 \\
Protected pocket faces & 474 & 474 \\
\bottomrule
\end{tabular}
\end{table}

The benchmark isolates this one CAD volume in air with the recorded GEMC3 executable and Geant4 stack,
FTFP\_BERT without optical processes, zero magnetic field, one thread, default production cuts, and no GUI
or event output. Both fixtures use the new volume name, identical physical placement attributes, and the
same alloy definition. Matched 5 GeV LUND primaries start at the target, uniform in solid angle over
0.25--3.0 degrees and full azimuth. Five alternating-order repetitions use seeds 12345--12349 and 200 events
per electron, proton, positive-pion, or neutron run. Separate zero-event runs measure startup without
overlap checking; time since beamOn measures transport. CPU time and peak RSS are retained separately.

Table~\ref{tab:inner-shield-timing} gives median transport wall seconds per 200 events.

\begin{table}[htbp]
\centering
\caption{Median transport wall seconds over five repetitions.}
\label{tab:inner-shield-timing}
\begin{tabular}{lrrr}
\toprule
Particle & Original & Optimized & Original/optimized \\
\midrule
e- & 18.565 & 15.048 & 1.234 \\
proton & 2.790 & 2.426 & 1.150 \\
pi+ & 3.655 & 3.292 & 1.110 \\
neutron & 2.504 & 2.425 & 1.033 \\
\bottomrule
\end{tabular}
\end{table}

Median startup simulation wall time changes from 0.9492 to 0.8435 s; peak loading RSS changes from 175.23 to
151.28 MiB. Median original/optimized transport wall-time ratios range from 1.033 to 1.234. All 52 processes
exit normally, and all forty transport runs report the 200 requested events. Results describe this isolated
workload and are not pooled with earlier beamline or LTCC benchmarks.

Table~\ref{tab:inner-shield-warnings} gives Geant4 exception counts over all 52 processes.

\begin{table}[htbp]
\centering
\caption{Geant4 exception counts for the isolated inner-shield samples.}
\label{tab:inner-shield-warnings}
\begin{tabular}{lrr}
\toprule
Diagnostic & Original & Optimized \\
\midrule
GeomSolids1001 & 26 & 0 \\
GeomNav1002 & 0 & 0 \\
GeomNav0003 & 0 & 0 \\
GeomVol1002 & 0 & 0 \\
\bottomrule
\end{tabular}
\end{table}

The diagnostic table records the original wrong-orientation import warning and all other Geant4 exceptions
in the sampled runs. Full default beamline construction completes without a missing mesh, with 4 solid
warnings and 11 overlap reports, compared with five and ten before this change. The full check reports no
overlap involving the optimized inner shield; the listed overlap reports concern other volumes. Their
sampled count can change when CAD loading order and surface tessellation change. All fifteen native
geometry/material variations and fourteen CAD variations pass the semantic comparison against GEMC2. Matched
primaries and seeds do not ensure identical stochastic shower histories, and completion does not establish
detector-response equivalence or complete overlap freedom.

\FloatBarrier
Figure~\ref{fig:inner-shield-stl} compares the stored triangulations.

\begin{figure}[htbp]
\centering
\includegraphics[width=\textwidth]{figures/inner_shield_stl_comparison.pdf}
\caption{Original (left) and optimized (right) FT inner shield and flange STL surfaces.
Matched cameras and bounds use original STL coordinates; the small enclosed pocket is retained exactly.}
\label{fig:inner-shield-stl}
\end{figure}

\FloatBarrier
\section{FT inner shield bolts}

The sixth beamline study renames nut1--nut9 to FTInnerShieldAndFlange\_Bolt1 through
FTInnerShieldAndFlange\_Bolt9. These engineering meshes are solid copper bolts, despite their former names.
The approved simplification removes threads and replaces the hexagonal head details with round heads and
smooth solid shafts. Each surface uses 24 angular sides, an annular shoulder, and closed end caps: 192
triangles per bolt. No hollow bore is introduced. G4\_Cu, colors, mothers, and passive placements are
retained across ten applicable CAD variations and ninety GEMC2 SQLite rows; only their names change in the
placement records.

The shaft axes, head shoulders, and outer ends are measured from each original mesh. The head radius is
14.2875 mm, within the original hexagonal across-flats envelope. The original tip caps contact the shield
blind-hole end faces. Shortening each inner tip by 0.15 mm creates an explicit gap; the smooth shaft radius
is then solved from the original signed solid volume. The seven regular shafts have radius approximately
8.5741 mm; the two longer shafts, bolts 4 and 8, have radius approximately 8.3681 mm. Their individual
float32 STL volumes agree with the originals within 0.0001\%. This redistributes copper locally and changes
external surfaces rather than preserving the engineering threaded profile.

Combined triangles fall from 16,278 to 1,728 (89.38\% fewer), and bytes fall from 814,656 to 87,156.
Original and optimized total copper volumes are 150623.172389 and 150623.196214 mm$^3$, a change of
0.00001582\%. All nine outputs are closed, singly connected, manifold genus-zero surfaces with no detected
self-intersections or zero-area triangles. The archived originals reproduce all nine optimized STLs byte for
byte.

Table~\ref{tab:bolts-meshes} records the individual smooth-shaft models.

\begin{table}[htbp]
\centering
\caption{Original and simplified copper bolt models; volume change is in percent.}
\label{tab:bolts-meshes}
\begin{tabular}{lrrrr}
\toprule
Bolt & Original faces & New faces & Shaft radius (mm) & Volume change (\%) \\
\midrule
1 & 1,690 & 192 & 8.5741 & 0.0000057 \\
2 & 1,690 & 192 & 8.5741 & 0.0000107 \\
3 & 1,690 & 192 & 8.5741 & 0.0000037 \\
4 & 2,224 & 192 & 8.3681 & 0.0000899 \\
5 & 1,690 & 192 & 8.5741 & 0.0000037 \\
6 & 1,690 & 192 & 8.5741 & 0.0000057 \\
7 & 1,690 & 192 & 8.5741 & 0.0000057 \\
8 & 2,224 & 192 & 8.3681 & 0.0000029 \\
9 & 1,690 & 192 & 8.5741 & 0.0000057 \\
\bottomrule
\end{tabular}
\end{table}

The optimized FTInnerShieldAndFlange STL is unchanged. Since shield and bolts share the same placement
transform, their collision test uses stored STL coordinates. VTK triangle collision finds zero surface
contacts, and enclosed-point checks in both directions find no bolt vertices inside the shield and no shield
vertices inside any bolt. Minimum bidirectional vertex-to-surface distances are 0.150002 mm for each bolt.
These sampled distances are not a certified global clearance bound; the independent surface and containment
checks address intersection. A ten-volume shield-plus-bolts Geant4 fixture completes without exceptions at
check\_overlaps=2 and again with a 10,000-point geometry test. The unchanged shield hash is retained with
the clearance results.

Tracking benchmarks isolate the nine copper bolts in air, excluding the shield from both timed fixtures.
Both cases use the new names, identical materials and placements, the same recorded GEMC3 executable,
FTFP\_BERT without optical processes, no magnetic field, one thread, default production cuts, and no GUI or
event output. Matched 5 GeV LUND primaries start at the target and sample solid angle uniformly from 1.3 to
3.0 degrees over full azimuth. Five repetitions with alternating fixture order and seeds 12345--12349 use
200 events per electron, proton, positive-pion, and neutron run. Time since beamOn measures transport;
separate zero-event runs measure startup. All 52 processes complete normally with all forty event counts
correct and no recorded Geant4 exceptions in either case.

Table~\ref{tab:bolts-timing} gives median transport wall seconds per 200 events.

\begin{table}[htbp]
\centering
\caption{Median transport wall seconds over five repetitions of the nine-bolt workload.}
\label{tab:bolts-timing}
\begin{tabular}{lrrr}
\toprule
Particle & Original & Optimized & Original/optimized \\
\midrule
e- & 0.551724 & 0.457388 & 1.206 \\
proton & 0.329016 & 0.329104 & 1.000 \\
pi+ & 0.330425 & 0.283490 & 1.166 \\
neutron & 0.133233 & 0.111502 & 1.195 \\
\bottomrule
\end{tabular}
\end{table}

Median startup simulation wall time is 0.659214 s originally and 0.623648 s after simplification. Peak
startup RSS falls from 129.34 to 115.31 MiB. The transport wall-time ratios are 1.206 for electrons, 1.000
for protons, 1.166 for positive pions, and 1.195 for neutrons. The proton result shows essentially unchanged
tracking time. These short workloads and five repetitions do not establish a universal speedup; material
redistribution changes stochastic shower histories even for matched primaries and seeds. Full CLAS12
performance and detector-response equivalence are not inferred.

Full default beamline construction completes without missing meshes and with no reported bolt-shield
overlap. Four remaining solid warnings and ten overlap reports concern other volumes, compared with four and
eleven immediately before this study. Sampled counts depend on tessellation and loading order, so this
difference is not claimed as a general overlap repair. All fifteen native geometry/material variations and
fourteen CAD variations pass the semantic comparator against GEMC2. Earlier apex, downstream,
adapter/exit-shield, adjuster, and inner-shield tracking measurements remain in their respective sections
with separate workloads and archives.

\FloatBarrier
Figure~\ref{fig:bolts-stl} compares a regular and a longer copper bolt.

\begin{figure}[htbp]
\centering
\includegraphics[width=\textwidth]{figures/bolts_stl_comparison.pdf}
\caption{Original threaded bolts (left) and smooth solid replacements (right).
Bolts 1 and 4 represent the regular and longer models; each pair shares bounds and camera settings.
Coordinates are recentered and permuted to put the bolt shaft on the horizontal axis.}
\label{fig:bolts-stl}
\end{figure}

\FloatBarrier
\section{FT outer flange}

The seventh beamline study renames outerFlange to FTOuterFlange in ten applicable source and runtime CAD
variations and ten GEMC2 SQLite placement rows. The stainless-steel material, color, mother, and passive
placement are retained, including the original 1272.77 mm z translation. That translation differs by 0.5 mm
from the inner-shield and outer-mount placements and is not rounded to their common value.

The original is a closed manifold annular plate with 2,132 triangles and genus 13. The authorized hole
removal closes twelve bolt tunnels with pygemc's remove\_holes utility, leaving the central opening intact
and reducing genus to 1. Flat annular faces are then retriangulated using every original inner and outer rim
coordinate on both z faces. The inner radius is 95.25 mm, the outer radius 127 mm, and the original
thickness approximately 19.049927 mm. Each face retains 97 inner rim vertices and 112 outer rim vertices.
Saved coordinates are verified exactly against those engineering vertices; the original bounds are
unchanged.

Triangles fall to 836 (60.79\% fewer), and bytes fall from 106,684 to 41,884. The output remains one closed
manifold surface with no detected self-intersections or zero-area faces. Closing the bolt holes changes
volume from 389837.402125 to 422180.789284 mm$^3$, an increase of 8.296635\%. Retriangulation changes the
hole-filled reference volume only at floating-point precision. Bidirectional vertex, edge, and surface
sampling against that reference gives a maximum difference of 0.0000148 mm. This sampled difference
describes retriangulation, not the deliberately closed holes, and is not a proven global error bound. The
archived engineering mesh reproduces the optimized STL byte for byte.

Table~\ref{tab:outer-flange-mesh} summarizes the original and hole-filled flange.

\begin{table}[htbp]
\centering
\caption{Original and optimized FT outer flange.}
\label{tab:outer-flange-mesh}
\begin{tabular}{lrr}
\toprule
Quantity & Original & Optimized \\
\midrule
Triangles & 2,132 & 836 \\
Bytes & 106,684 & 41,884 \\
Genus & 13 & 1 \\
Volume (mm$^3$) & 389837.402 & 422180.789 \\
\bottomrule
\end{tabular}
\end{table}

The isolated tracking fixture contains only this stainless-steel flange in air. Both cases use the new name
and identical material and placement attributes. Matched 5 GeV LUND primaries originate at the target,
uniform in solid angle over 1.9--2.8 degrees and full azimuth. FTFP\_BERT runs without optical processes or
magnetic field, with one thread, default production cuts, and no GUI or event output. Five alternating-order
repetitions use seeds 12345--12349 and 200 events for each electron, proton, positive-pion, and neutron run.
Separate zero-event runs measure startup; elapsed time since beamOn measures transport. A fixed executable
snapshot is used for every run, with its hash checked afterward. Earlier runs interrupted by sandbox
resource-timer restrictions or invalidated by a changing shared executable are excluded.

Table~\ref{tab:outer-flange-timing} records the tracking measurements.

\begin{table}[htbp]
\centering
\caption{Median transport wall seconds per 200 events over five repetitions.}
\label{tab:outer-flange-timing}
\begin{tabular}{lrrr}
\toprule
Particle & Original & Optimized & Original/optimized \\
\midrule
e- & 0.646804 & 0.668927 & 0.967 \\
proton & 0.429432 & 0.425280 & 1.010 \\
pi+ & 0.393231 & 0.377190 & 1.043 \\
neutron & 0.221017 & 0.274356 & 0.806 \\
\bottomrule
\end{tabular}
\end{table}

Median startup simulation wall time changes from 0.699890 to 0.693861 s; peak startup RSS changes from
117.83 to 115.67 MiB. Median original/optimized transport ratios are 0.967, 1.010, 1.043, and 0.806 for
electrons, protons, positive pions, and neutrons, respectively. A ratio above one means shorter optimized
tracking time; a ratio below one means longer time. All 52 processes complete normally, and all forty
transport runs report their 200 requested events. These short isolated workloads do not establish a
universal tracking gain. The added stainless steel changes scattering and stochastic shower histories, so
timing differences cannot be attributed solely to triangle reduction or treated as detector-response
equivalence.

Table~\ref{tab:outer-flange-warnings} records isolated-fixture exceptions across all 52 processes.

\begin{table}[htbp]
\centering
\caption{Geant4 exception counts in the timed isolated flange fixtures.}
\label{tab:outer-flange-warnings}
\begin{tabular}{lrr}
\toprule
Diagnostic & Original & Optimized \\
\midrule
GeomSolids1001 & 0 & 0 \\
GeomNav1002 & 0 & 0 \\
GeomNav0003 & 0 & 0 \\
GeomVol1002 & 0 & 0 \\
\bottomrule
\end{tabular}
\end{table}

A separate twelve-volume assembly fixture includes the flange, unchanged inner shield and outer mount, and
nine copper bolts. Both original and optimized fixtures use the new flange name and the same placements.
Their 10,000-point Geant4 checks report 5 and 5 overlap warnings, respectively; both also retain one
surface-hole import warning from the unchanged outer mount. The retained logs identify the volumes and
micrometre-scale penetrations. Existing outer-mount/flange interference remains a limitation of the
assembly, so no claim of complete overlap freedom is made. Full default construction completes without
missing meshes, with 4 solid warnings and 9 overlap reports. All fifteen native geometry/material variations
and fourteen CAD variations pass the semantic comparison against GEMC2. The earlier six beamline studies
retain their independent timing tables, diagnostic results, and archives.

The subsequent mount study resolves this earlier joint interference; see Section~\ref{sec:outer-mount}.

\subsection{Tracking at fixed material volume}

A follow-up control compares the hole-filled flange before and after retriangulation, removing the added
material as a difference between the cases. The baseline has 2,084 triangles and 104,284 bytes; the
simplified flange has 836 triangles and 41,884 bytes, a 59.88\% triangle reduction relative to this filled
reference. Both have the same central opening, bounds, stainless-steel material, placement, and volume of
422180.789284 cubic millimetres to the reported precision. Their sampled surface differences remain below
0.0000148 mm. The fixture label original now denotes the unreduced hole-filled reference, not the drilled
engineering mesh.

The control uses the same fixed executable snapshot and physics settings as the earlier flange study, with 5
GeV primaries uniform in solid angle over 1.9--2.8 degrees and full azimuth. Five alternating-order
repetitions use seeds 12345--12349. Each particle run has 2,000 events, ten times the earlier sample; both
cases have identical primary hashes and geometry/material database records. Separate startup samples retain
the loading and peak-memory measurements. The larger control is reported independently, without pooling its
absolute times or ratios with the earlier 200-event drilled-versus-filled results.

Table~\ref{tab:outer-flange-control} gives median transport wall seconds per 2,000 events.

\begin{table}[htbp]
\centering
\caption{Hole-filled reference versus simplified flange: medians over five repetitions.}
\label{tab:outer-flange-control}
\begin{tabular}{lrrr}
\toprule
Particle & Filled reference & Simplified & Reference/simplified \\
\midrule
e- & 7.40750 & 6.90088 & 1.073 \\
proton & 4.63719 & 4.56330 & 1.016 \\
pi+ & 4.19454 & 4.05502 & 1.034 \\
neutron & 2.94839 & 2.91016 & 1.013 \\
\bottomrule
\end{tabular}
\end{table}

Median baseline/simplified tracking wall-time ratios are 1.073 for electrons, 1.016 for protons, 1.034 for
positive pions, and 1.013 for neutrons. A value above one means shorter tracking time for the simplified
mesh. These measurements quantify retriangulation at fixed material volume. The simplified mesh is faster in
all twenty particle/seed pairs; the individual paired ratios are retained in the control dataset. Matched
primaries and seeds do not guarantee identical subsequent shower histories because tessellated boundary
arithmetic can alter stochastic transport. Five repetitions do not establish a universal speedup or a
precise causal decomposition of the earlier slowdown.

Median startup simulation wall time changes from 0.704797 to 0.692080 s, with peak startup RSS changing from
117.52 to 115.58 MiB. All 52 control processes complete normally, and all forty transport runs report the
requested 2,000 events. Geant4 exception counts, paired seed ratios, fixture hashes, and raw timing reports
are retained in a separate control dataset. Neither case records a Geant4 exception in the sampled control.

\FloatBarrier
Figure~\ref{fig:outer-flange-stl} compares the original and filled triangulations.

\begin{figure}[htbp]
\centering
\includegraphics[width=\textwidth]{figures/outer_flange_stl_comparison.pdf}
\caption{Original FT outer flange (left) and optimized flange with twelve bolt holes filled (right).
Matched cameras and bounds use original STL coordinates; both retain the central opening and rim vertices.}
\label{fig:outer-flange-stl}
\end{figure}

\FloatBarrier
\section{FT outer mount}
\label{sec:outer-mount}

The eighth beamline study renames outerMount to FTOuterMount across ten source and runtime CAD variations
and ten GEMC2 SQLite rows. The stainless-steel material, color, mother, and passive placement are retained.
The mount remains at z = 1273.27 mm and the unchanged FTOuterFlange at z = 1272.77 mm, both with the
original passive 180-degree y rotation. Their 0.5 mm relative translation is included in every clearance
calculation.

The original mount contains 3,094 triangles, twelve boundary edges forming four triangular surface gaps, and
one manifold surface component. Its open surface has no defined solid volume. Closing these four missing
faces gives a 3,098-triangle sealed reference of genus 11 and volume 3889831.731170 mm$^3$. The central
passage and existing drilled openings are retained. This repaired reference defines the reported
material-volume comparison; the original engineering solid mass is not independently certified.

The original mount/flange check reports small interference at their nominal 95.25 mm radius because the two
cylindrical surfaces have different polygonal sampling. The flange inner polygon has minimum inradius
95.200045 mm. The mount lower outer rim is reduced to 95.150045 mm, approximately 0.100 mm of radial relief.
Adjacent sidewall facets interpolate toward unchanged upper vertices. All 883 facets adjoining the relieved
end are protected at their saved float32 coordinates during topology-preserving quadric simplification. The
flange STL is unchanged, and no placement is shifted to conceal the interference.

The final mount contains 2,740 triangles, 11.44\% fewer than the raw engineering mesh, and 137,084 bytes
rather than 154,784. It is one closed manifold component of genus 11 with no detected self-intersections or
zero-area triangles. Bidirectional sampling against the relieved reference gives a maximum difference of
0.042078 mm, below the 0.05 mm guard. These are sampled differences rather than proven global error bounds.
Volume changes by -0.249235\% against the sealed reference; simplification alone changes the relieved
reference volume by 0.00011056\%. The axial bounds are retained, with at most about 0.00065 mm change in the
transverse bounding coordinates. The builder reproduces the final STL byte for byte from the archived
original.

Table~\ref{tab:outer-mount-mesh} records the original and optimized mount inventory.

\begin{table}[htbp]
\centering
\caption{Original and optimized FT outer mount.}
\label{tab:outer-mount-mesh}
\begin{tabular}{lrr}
\toprule
Quantity & Original & Optimized \\
\midrule
Triangles & 3,094 & 2,740 \\
Bytes & 154,784 & 137,084 \\
Boundary edges & 12 & 0 \\
Surface components & 1 & 1 \\
Solid volume (mm$^3$) & Undefined & 3880136.906 \\
\bottomrule
\end{tabular}
\end{table}

Placed-mesh verification finds zero triangle contacts and no contained vertices in either direction between
mount and flange. In mount coordinates, the flange is shifted by +0.5 mm along z. Each mount triangle is
clipped to the flange z interval, and convexity of the radial norm bounds its radius by the clipped polygon
vertices. The maximum is 95.164071 mm, giving at least 0.035974 mm radial separation from the flange inner
polygon throughout their common axial interval. The minimum bidirectional vertex-to-surface distance is
0.085929 mm; this latter value is sampled and is not presented as a global surface-clearance bound. All ten
placement variants retain the same verified relative transform.

A twelve-volume assembly fixture includes the mount, unchanged inner shield and flange, and nine copper
bolts. With both original and optimized fixtures using the new mount name and identical placements, a
10,000-point Geant4 test gives five flange/mount overlap reports and one surface-hole warning for the
original mount. The optimized assembly completes with no Geant4 exceptions or overlap reports, resolving the
joint interference documented in the preceding flange study. Full default beamline construction also
completes without missing meshes or a mount/flange overlap. Its three solid warnings and eight overlap
reports concern other geometry, so this study does not establish complete beamline overlap freedom. All
fifteen native geometry/material variations and fourteen CAD variations pass the semantic comparison against
GEMC2.

The tracking benchmark isolates this one mount in air. Both cases use the new name, the same material and
placement, and the unchanged fixed GEMC3 executable snapshot from the flange study. Matched 5 GeV LUND
primaries begin at the target and sample solid angle uniformly over 1.5--2.4 degrees and full azimuth.
FTFP\_BERT runs without optical processes or magnetic field, using one thread, default production cuts, and
no GUI or event output. Five alternating-order repetitions use seeds 12345--12349 and 200 events per
electron, proton, positive-pion, and neutron run. Elapsed time since beamOn measures transport; separate
zero-event samples measure startup and peak RSS.

Table~\ref{tab:outer-mount-timing} records median transport wall seconds per 200 events.

\begin{table}[htbp]
\centering
\caption{Median transport wall seconds over five repetitions.}
\label{tab:outer-mount-timing}
\begin{tabular}{lrrr}
\toprule
Particle & Original & Optimized & Original/optimized \\
\midrule
e- & 15.97580 & 15.51250 & 1.030 \\
proton & 2.13591 & 2.10357 & 1.015 \\
pi+ & 2.93950 & 2.85889 & 1.028 \\
neutron & 1.83088 & 1.93090 & 0.948 \\
\bottomrule
\end{tabular}
\end{table}

Median original/optimized transport wall-time ratios are 1.030 for electrons, 1.015 for protons, 1.028 for
positive pions, and 0.948 for neutrons. A ratio below one denotes longer optimized tracking time. All 52
processes complete normally, and all forty transport runs report the requested 200 events. Startup
simulation wall time changes from 0.705080 to 0.693238 s, and peak startup RSS changes from 120.12 to 119.38
MiB. Closure repairs and relief change material boundaries, so these short workload timings are not
attributed solely to triangle reduction. They do not establish detector-response equivalence or full CLAS12
speedup. The preceding flange control at fixed material volume remains separately reported.

Table~\ref{tab:outer-mount-warnings} gives isolated-fixture exceptions across all 52 processes.

\begin{table}[htbp]
\centering
\caption{Geant4 exception counts in the isolated mount fixtures.}
\label{tab:outer-mount-warnings}
\begin{tabular}{lrr}
\toprule
Diagnostic & Original & Optimized \\
\midrule
GeomSolids1001 & 26 & 0 \\
GeomNav1002 & 0 & 0 \\
GeomNav0003 & 0 & 0 \\
GeomVol1002 & 0 & 0 \\
\bottomrule
\end{tabular}
\end{table}

\FloatBarrier
Figure~\ref{fig:outer-mount-stl} compares the stored mount triangulations.

\begin{figure}[htbp]
\centering
\includegraphics[width=\textwidth]{figures/outer_mount_stl_comparison.pdf}
\caption{Original FT outer mount (left) and repaired, relieved, and simplified mount (right).
Matched cameras and bounds use original STL coordinates; existing drilled openings are retained.}
\label{fig:outer-mount-stl}
\end{figure}

\FloatBarrier
\section{FT inner shield}
\label{sec:ft-shield}

The shield/cone interference recorded below is resolved by the subsequent cone revision in
Section~\ref{sec:moller-cone}.

The initial revision below is retained with its original measurements. The final smooth-flange configuration
supersedes it; see Section~\ref{sec:ft-shield-smooth}.

The ninth beamline study renames taggerInnerShield to FTInnerShield in all nine source and runtime CAD
variations that contain it and in nine GEMC2 SQLite CAD rows. The tungsten-alloy material beamline\_W,
color, mother, and original passive placement are retained: z = 1643.00 mm with a 180-degree y rotation.
This downstream grooved tube is distinct from the previously optimized FTInnerShieldAndFlange at z = 1273.27
mm.

The original contains 44,122 triangles, 7,404 boundary edges, 67 surface components, 351 non-manifold
vertices, and 1,750 faces selected by the intersection check. Its open and non-manifold surface has no
defined solid volume. Routine welding and generic hole filling do not establish a sound surface. A
deterministic reconstruction uses the measured tube cross-section, chamfers, flange fillets, and feature
positions, without Poisson reconstruction. The central bore remains 60 mm in diameter, nominal tube outer
diameter 80 mm, length 665 mm, and flange outer radius 100 mm. The flange occupies local z = -625 to -605
mm.

Two longitudinal grooves at local azimuths 45 and 135 degrees are retained after the flange. Each is 6 mm
wide and 3 mm deep, extending from local z = -575 mm to the tube end at z = 0. A 6 mm axial end radius
blends its floor into the tube. All 26 flange through-holes and the single radial tube hole are closed. The
requested 0.5 mm recess threshold is applied to axial depth: eight 5.75 mm deep conical recesses, sixteen
0.6 mm face chamfers, and two 10 mm deep counterbores remain as blind features. The conical mouths
originally touch the flange outer rim exactly; reducing each mouth radius from 10.00 to 9.95 mm removes that
tangency while retaining their depths and 4.25 mm bottom radii.

The reference uses 128 angular segments for the main profile, sampled fillet/end curves, and Boolean cuts
for the grooves and retained recesses. A 0.000001 mm weld removes near-coincident Boolean vertices and their
sliver facets before saving the reference. Its 15,100 triangles define a closed manifold component of genus
one with volume 1938424.030371 mm$^3$. Protected quadric simplification retains 9,404 groove, recess, and
corner faces at their saved float32 coordinates and produces 13,500 triangles, a 69.40\% reduction from the
engineering mesh. File size decreases from 2,206,184 to 675,084 bytes.

The final mesh remains one closed manifold component of genus one, with no selected self-intersections or
zero-area facets. Bounds remain [-100, -100, -665] to [100, 100, 0] mm apart from negligible source z
rounding. Minimum triangle altitude is 0.000344255 mm. Bidirectional sampling against the rebuilt reference
gives a maximum change of 0.023925 mm, below the 0.05 mm guard. Simplification changes reconstructed volume
by -0.00058479\%. These are sampled checks, not proven global surface bounds or a comparison to a certified
original mass.

Independent VTK containment probes verify all 58 requested inside/outside results across groove floors,
blind recesses, filled bolt shafts, the filled tube hole, and the central passage. Outside the removed
drilling walls and flange drilling region, retained source profile vertices lie within 0.050719 mm of the
reconstructed surface. One-way vertex, edge, and face sampling of the retained raw tube/rim subset gives a
larger maximum 0.969091 mm: the reconstruction replaces coarse, disconnected raw facets with the measured
nominal profile. This original-to-final difference is distinct from the much smaller subsequent decimation
error. The original mass change and detector-response equivalence are not established. The builder
reproduces the final STL and reconstructed reference byte for byte from the archived engineering input.

\begin{table}[htbp]
\centering
\caption{Original and optimized downstream FT inner shield.}
\label{tab:ft-shield-mesh}
\begin{tabular}{lrr}
\toprule
Quantity & Original & Optimized \\
\midrule
Triangles & 44,122 & 13,500 \\
Bytes & 2,206,184 & 675,084 \\
Boundary edges & 7,404 & 0 \\
Surface components & 67 & 1 \\
Solid volume (mm$^3$) & Undefined & 1938412.695 \\
\bottomrule
\end{tabular}
\end{table}

A two-volume assembly includes the unchanged main-cone at z = 1589.27 mm and the shield at its original
placement. The 10,000-point Geant4 test reports five shield/cone overlaps in each case, confirming the
pre-existing interference remains. Original construction has three solid warnings; optimized assembly
construction has one warning from the unchanged cone. The isolated optimized shield loads with no Geant4
exceptions. Full default beamline construction completes without missing meshes, with 1 solid warning and 9
overlap reports involving the shield/cone and downstream/adapter geometry. The original placements are
retained rather than adjusted in this study. All fifteen native geometry/material and fourteen CAD
variations pass the semantic comparison against GEMC2.

Two independent tracking comparisons isolate the shield in air. The first uses the defective engineering
mesh as baseline; the second uses the hole-filled, grooved, unreduced reference to measure simplification
after repair. Each uses the fixed GEMC3 executable snapshot from the flange study, matched 5 GeV
target-origin primaries uniform in solid angle over 0.7--2.7 degrees and full azimuth, FTFP\_BERT without
optics or magnetic field, one thread, default production cuts, and no GUI or event output. Five
alternating-order repetitions use seeds 12345--12349 and 200 events per electron, proton, positive-pion, and
neutron run. Each comparison contains 52 processes, including zero-event startup and isolated overlap
checks. All 104 processes complete normally and all eighty transport runs report 200 events.

Engineering-original/optimized median transport wall-time ratios are 0.923, 1.007, 0.992, 0.924 for
electrons, protons, positive pions, and neutrons respectively. Ratios below one denote slower optimized
tracking. The engineering surface carries solid and navigation warnings; repair and authorized hole closure
change material boundaries, so this comparison does not attribute timing changes solely to triangle count.
Original material volume is undefined.

\begin{table}[htbp]
\centering
\caption{Median transport wall seconds per 200 events: engineering comparison.}
\label{tab:ft-shield-engineering}
\begin{tabular}{lrrr}
\toprule
Particle & Original & Optimized & Original/optimized \\
\midrule
e- & 22.46020 & 24.33720 & 0.923 \\
proton & 2.23604 & 2.22142 & 1.007 \\
pi+ & 2.96617 & 2.99009 & 0.992 \\
neutron & 1.99370 & 2.15824 & 0.924 \\
\bottomrule
\end{tabular}
\end{table}

Rebuilt-reference/optimized median ratios are 1.087, 1.075, 1.075, 1.049 in the same particle order. This
comparison retains the same blind recesses, closed holes, grooves, and material, with the recorded small
decimation volume difference. Both isolated reference and optimized fixtures have no Geant4 exceptions.
These short workload results are reported separately from the repair comparison; small ratios are sensitive
to timing variation and stochastic transport. They do not establish full CLAS12 speedup or response
equivalence. The optimized mesh is faster in all 20 matched particle/seed pairs in the reference comparison.
Median startup simulation time changes from 1.063410 to 0.770623 s, with peak startup RSS 199.94 to 136.06
MiB.

\begin{table}[htbp]
\centering
\caption{Median transport wall seconds per 200 events: reconstructed-reference comparison.}
\label{tab:ft-shield-reference}
\begin{tabular}{lrrr}
\toprule
Particle & Reference & Optimized & Reference/optimized \\
\midrule
e- & 26.31760 & 24.21180 & 1.087 \\
proton & 2.37224 & 2.20600 & 1.075 \\
pi+ & 3.16977 & 2.94966 & 1.075 \\
neutron & 2.28295 & 2.17616 & 1.049 \\
\bottomrule
\end{tabular}
\end{table}

\begin{table}[htbp]
\centering
\caption{Geant4 exception counts across each isolated 52-process comparison.}
\label{tab:ft-shield-warnings}
\begin{tabular}{lrrr}
\toprule
Diagnostic & Original & Optimized & Reference \\
\midrule
GeomSolids1001 & 52 & 0 & 0 \\
GeomSolids1002 & 0 & 0 & 0 \\
GeomNav1002 & 45 & 0 & 0 \\
GeomNav0003 & 0 & 0 & 0 \\
GeomVol1002 & 0 & 0 & 0 \\
\bottomrule
\end{tabular}
\end{table}

\FloatBarrier
\begin{figure}[htbp]
\centering
\includegraphics[width=0.9\textwidth]{figures/ft_shield_stl_comparison.pdf}
\caption{Original FT inner shield (left) and rebuilt, simplified shield (right). Matched cameras and bounds
use original coordinates; longitudinal grooves are retained.}
\label{fig:ft-shield-stl-comparison}
\end{figure}

\FloatBarrier
\begin{figure}[htbp]
\centering
\includegraphics[width=0.9\textwidth]{figures/ft_shield_details.pdf}
\caption{Matched flange underside views and tube sections at local z = -100 mm. Flange recesses remain
blind, and both longitudinal grooves remain after through-hole closure.}
\label{fig:ft-shield-details}
\end{figure}

\FloatBarrier
\subsection{Smooth flange revision}
\label{sec:ft-shield-smooth}

The shield/cone interference recorded below is resolved by the subsequent cone revision in
Section~\ref{sec:moller-cone}.

The smooth-flange FTInnerShield revision fills eight conical recesses, sixteen face chamfers, and two
stepped counterbores retained in the previous 13,500-triangle version. Both flange faces are complete flat
annuli. The large central passage, tube chamfers and flange fillets, and both longitudinal tube grooves
remain. The name, tungsten material, color, and all nine passive placements remain unchanged.

Removing the recess cuts reduces the reconstructed reference to 8,536 triangles. Protected simplification
gives 6,500 triangles, 51.85\% fewer than the previous revision and 85.27\% fewer than the defective
44,122-triangle engineering input. File size decreases from 675,084 to 325,084 bytes. All 3,148 groove,
flat-flange, and corner faces are protected, including all 512 flange-plane facets. Further trials at 6,300
and 6,200 triangles exceed the shape/quality guard and are not used. The result is one closed manifold
component of genus one with no detected self-intersections or zero-area facets, and unchanged bounds.

The previous closed volume is 1938412.694702 mm$^3$ and the final volume is 1946864.540816 mm$^3$. Filling
the recesses and applying the new simplification changes tungsten volume by +0.436019\%. The smooth,
unreduced reference has volume 1946881.004462 mm$^3$; simplification alone changes it by -0.00084564\%.
Bidirectional sampling against that reference gives maximum displacement 0.047853 mm, below the 0.05 mm
guard. These sampled limits apply after the authorized recess closure, rather than across the filled
cavities. The original engineering mass remains undefined.

Independent surface checks confirm 256 coplanar triangles at each flange face, with areas matching the
complete 128-sided annuli at local z = -625 and -605 mm. The inner face radii are 45 and 50 mm respectively,
adjoining the retained fillets; outer radius is 100 mm. Every outer flange surface facet belongs to one flat
face or the cylindrical rim, with zero recess-wall facets. All 1,134 independent containment probes pass,
including a dense grid just inside both flange faces, all former bolt and recess centres, tube groove
floors, the filled tube hole, and the central beam passage. The final STL and smooth reference reproduce
byte for byte from the archived raw input.

\begin{table}[htbp]
\centering
\caption{FT inner-shield smooth-flange revision.}
\label{tab:ft-shield-smooth-mesh}
\begin{tabular}{lrr}
\toprule
Quantity & Previous & Smooth \\
\midrule
Triangles & 13,500 & 6,500 \\
Bytes & 675,084 & 325,084 \\
Flange recesses & 26 & 0 \\
Boundary edges & 0 & 0 \\
Volume (mm$^3$) & 1938412.695 & 1946864.541 \\
\bottomrule
\end{tabular}
\end{table}

Two separate tracking comparisons use the previous closed 13,500-triangle mesh and the new smooth, unreduced
reference as baselines respectively. Both use the same fixed GEMC3 executable snapshot, beamline\_W
material, original placement, and matched 5 GeV target-origin primaries over 0.7--2.7 degrees and full
azimuth. FTFP\_BERT runs without optics or magnetic field, with one thread, default production cuts, and no
GUI or event output. Five alternating-order repetitions use seeds 12345--12349 and 200 events per electron,
proton, positive-pion, and neutron run. Each comparison has 52 processes including startup and isolated
overlap checks; all 104 complete normally and all eighty transport runs report 200 events, with no Geant4
exceptions in any fixture.

Previous/smooth median transport ratios are 1.273, 1.152, 1.160, 1.134 for electrons, protons, positive
pions, and neutrons respectively. These incremental upgrade timings include the added tungsten from filling
the blind recesses. Reference/smooth medians are 1.055, 1.023, 1.041, 1.030 in the same particle order,
comparing both cases with fully smooth flange surfaces and the recorded small decimation volume difference.
A ratio above one denotes faster tracking. Both comparisons are retained separately; neither establishes
full CLAS12 speedup or detector-response equivalence. Median startup simulation time changes from 0.776447
to 0.739538 s, with peak startup RSS 135.95 to 124.98 MiB.

\begin{table}[htbp]
\centering
\caption{Median transport wall seconds per 200 events: incremental upgrade.}
\label{tab:ft-shield-smooth-upgrade}
\begin{tabular}{lrrr}
\toprule
Particle & Previous & Smooth & Previous/smooth \\
\midrule
e- & 23.96760 & 18.82900 & 1.273 \\
proton & 2.23387 & 1.93947 & 1.152 \\
pi+ & 2.97310 & 2.56412 & 1.160 \\
neutron & 2.18361 & 1.92635 & 1.134 \\
\bottomrule
\end{tabular}
\end{table}

\begin{table}[htbp]
\centering
\caption{Median transport wall seconds per 200 events: smooth-reference control.}
\label{tab:ft-shield-smooth-reference}
\begin{tabular}{lrrr}
\toprule
Particle & Reference & Smooth & Reference/smooth \\
\midrule
e- & 20.41750 & 19.34670 & 1.055 \\
proton & 1.98474 & 1.94027 & 1.023 \\
pi+ & 2.68023 & 2.57448 & 1.041 \\
neutron & 1.96960 & 1.91241 & 1.030 \\
\bottomrule
\end{tabular}
\end{table}

The unchanged main-cone is included in a separate two-volume 10,000-point assembly check. Both previous and
smooth cases retain five pre-existing shield/cone overlap reports and one solid warning from the cone. Full
default construction completes without missing meshes, with one solid warning and eight overlap reports in
other or pre-existing joints. Smooth flange closure does not resolve that separate cone interference. All
fifteen native geometry/material and fourteen CAD variations continue to pass semantic comparison against
GEMC2. The earlier study archive is retained unchanged; this revision has a separate archive with its
previous mesh, smooth reference, final mesh, complete raw benchmarks, and validation evidence.

\FloatBarrier
\begin{figure}[htbp]
\centering
\includegraphics[width=0.9\textwidth]{figures/ft_shield_smooth_comparison.pdf}
\caption{Previous flange with blind recesses (left) and final smooth flange (right), with matched tube
sections showing both longitudinal grooves and the central bore retained.}
\label{fig:ft-shield-smooth}
\end{figure}

\FloatBarrier
\section{Moller cone}
\label{sec:moller-cone}

The remaining downstream contacts recorded in this historical study are resolved by the subsequent socket
revision in Section~\ref{sec:downstream-joints}.

The tenth beamline study renames main-cone to moller\_cone in nine source and runtime GXML variations and
nine GEMC2 SQLite CAD rows. The beamline\_W tungsten material, color, passive 180-degree y rotation, and
run-dependent positions remain unchanged. Five placements use cone z = 1589.27 mm and overlap the
FTInnerShield axial interval; four use z = 1273.27 mm and are already separated from the shield. The current
smooth-flange FTInnerShield is unchanged. The raw engineering cone is archived before replacement.

The original contains 15,782 triangles and 302 boundary edges in eight surface loops: five on the outer
tapered wall and three on the inner cylindrical wall. Its open surface has no valid enclosed material
volume. Filling those openings follows the measured outer cone and inner cylinder profiles, preserving every
original face and vertex before clearance work. The closed hole-filled reference has 16,688 triangles. The
large central beam passage remains open. A full nominal axisymmetric reconstruction is avoided because the
engineering outer rings are coarsely tessellated; their existing bounds and facets provide the shape
reference.

The existing entrance counterbore, about 39.98595 mm radius, is widened with a protected 128-sided
cylindrical cut of radius 40.200001 mm. The cut ends at local cone z = -53.53 mm, 0.2 mm beyond the
overlapping shield tip at -53.73 mm. A 0.2 mm nominal radial allowance above the shield radius accounts for
the polygonal opening. Independent clipped-triangle bounds give a minimum cone opening radius of 40.187892
mm and maximum shield radius 40.000001 mm, proving at least 0.187891 mm radial separation throughout the
shared axial interval. Collision checks and bidirectional vertex containment find no intersection in all
nine placements; the nearest sampled vertex-to-surface distance is 0.199938 mm in the overlapping group and
199.109193 mm in the axially separated group.

The clearance reference has 11,616 triangles; protected simplification retains 1,964 joint, shoulder, and
corner facets and produces 5,000 triangles, 68.32\% fewer than the engineering input. File size falls from
789,184 to 250,084 bytes. The final mesh is one closed manifold component of genus one, with no detected
self-intersections or zero-area facets and exactly unchanged outer bounds. Minimum triangle altitude is
0.000380430 mm. Maximum sampled bidirectional displacement from the clearance reference is 0.020048 mm,
below the 0.05 mm shape guard. The original input reproduces the final STL and both reference meshes byte
for byte.

Hole-filled reference volume is 5835559.232219 mm$^3$; final volume is 5825708.437762 mm$^3$. Counterbore
relief and simplification remove 0.168806\% tungsten relative to the closed hole-filled reference.
Simplification alone changes clearance-reference volume by -0.0000002660\%. No original mass difference is
inferred from the defective open input. Fifty-five independent material probes test the filled opening
centres and the retained central beam passage.

\begin{table}[htbp]
\centering
\caption{Original, repaired clearance reference, and final Moller cone.}
\label{tab:moller-mesh}
\begin{tabular}{lrrr}
\toprule
Quantity & Original & Reference & Final \\
\midrule
Triangles & 15,782 & 11,616 & 5,000 \\
Bytes & 789,184 & 580,884 & 250,084 \\
Boundary edges & 302 & 0 & 0 \\
Volume (mm$^3$) & undefined & 5825708.453 & 5825708.438 \\
\bottomrule
\end{tabular}
\end{table}

Two tracking comparisons isolate the cone in air: raw engineering input versus final mesh and repaired
clearance reference versus final mesh. They use the same fixed GEMC3 binary, beamline\_W material, default
placement, and matched target-origin 5 GeV primaries over 1.1--2.6 degrees and full azimuth. FTFP\_BERT runs
without optics or field, with one thread, default production cuts, and no GUI or event output. Each
comparison uses five alternating order repetitions, seeds 12345--12349, and 200 events per electron, proton,
positive-pion, and neutron run. Each has 52 processes including startup and isolated overlap checks. All 104
processes complete normally and all eighty transport runs report 200 events.

Raw/final median transport ratios are 2.200, 1.739, 1.574, 1.552 for electrons, protons, positive pions, and
neutrons respectively. The raw cone emits 26 GeomSolids1001 and 396 GeomNav1002 exception blocks across the
recorded processes. Its stuck-track recovery and invalid solid prevent attributing these ratios solely to
tessellation reduction. The final mesh emits no Geant4 exceptions. The clean repaired-reference/final
comparison gives ratios 1.638, 1.313, 1.402, 1.301 in the same order, with no exceptions in either case.
Ratios above one mean faster tracking. These isolated timings do not establish full CLAS12 speedup or
physics-response equivalence. Median startup simulation time changes from 0.777853 to 0.744276 s; peak
startup RSS changes from 137.08 to 122.14 MiB.

\begin{table}[htbp]
\centering
\caption{Median transport wall seconds per 200 events: raw input versus final cone.}
\label{tab:moller-upgrade}
\begin{tabular}{lrrr}
\toprule
Particle & Original & Final & Original/final \\
\midrule
e- & 55.95160 & 25.43760 & 2.200 \\
proton & 10.28140 & 5.91253 & 1.739 \\
pi+ & 15.50770 & 9.85140 & 1.574 \\
neutron & 11.84350 & 7.63252 & 1.552 \\
\bottomrule
\end{tabular}
\end{table}

\begin{table}[htbp]
\centering
\caption{Median transport wall seconds per 200 events: repaired clearance-reference control.}
\label{tab:moller-reference}
\begin{tabular}{lrrr}
\toprule
Particle & Reference & Final & Reference/final \\
\midrule
e- & 28.23070 & 17.23620 & 1.638 \\
proton & 5.43417 & 4.13941 & 1.313 \\
pi+ & 8.36232 & 5.96464 & 1.402 \\
neutron & 5.71081 & 4.39053 & 1.301 \\
\bottomrule
\end{tabular}
\end{table}

A separate cone/shield assembly test uses 10,000-point Geant4 overlap checking. The original cone gives five
overlap reports and one solid warning; the final assembly has no warnings or overlaps. Full default beamline
construction completes with no missing meshes or solid warnings, retaining four overlap reports at the
unrelated downstream tube/adapter joints. The earlier shield studies retain their historical overlap
measurements; this cone revision resolves that interference. Semantic comparisons pass all fifteen native
geometry and material variations and fourteen CAD variations against GEMC2.

\FloatBarrier
\begin{figure}[htbp]
\centering
\includegraphics[width=0.9\textwidth]{figures/moller_cone_stl_comparison.pdf}
\caption{Original engineering cone (left) and optimized moller\_cone (right), with matched cameras and
bounds. Surface openings are filled and the shield counterbore is widened; the central beam passage and
outer envelope remain.}
\label{fig:moller-cone}
\end{figure}

\begin{figure}[htbp]
\centering
\includegraphics[width=0.95\textwidth]{figures/moller_cone_joint.pdf}
\caption{Actual triangle sections at y = 0 through the original and final cone with the unchanged FT inner
shield. The enlarged counterbore clears the shield tip; the right panel magnifies the former shoulder
interference. The guaranteed radial gap includes polygon faceting.}
\label{fig:moller-joint}
\end{figure}

\FloatBarrier
\section{Downstream joint clearance revision}
\label{sec:downstream-joints}

A subsequent revision resolves the remaining contacts between Downstream\_Beam\_tube\_Segment\_1 and both
Beamline\_Adapter\_1 and Downstream\_Beam\_tube\_Segment\_2. The nominal cylindrical gaps are 0.0508 and
0.0127 mm, but the mating surfaces use 24-sided polygons with different angular alignments. Both alignments
are present in the engineering inputs. The reported 0.0398313 and 0.0568888 mm Geant4 penetration samples
are not global maxima. Polygon support planes give deficits of 0.342860 and 0.395184 mm at the respective
joints. Independent triangle checks find 240 and 239 contacts in the previous meshes; positive
vertex-distance samples alone would not establish clearance.

Only Segment 1 is changed. Its entrance socket radius increases from 60.325 to 60.771702 mm through nominal
z = 5048.376367 mm, followed by a 1 mm transition back to the original bore. Two additional 24-vertex rings
replace one long wall band with three bands, adding 96 triangles. The downstream socket radius increases
from 70.6628 to 71.162277 mm over its existing downstream wall interval. Both sockets retain their angular
grids. The adapter, Segment 2, all outer bounds, materials, passive placements, and the open beam passage
remain unchanged.

The revised Segment 1 contains 5,096 triangles and 254,884 bytes, compared with 5,000 triangles and 250,084
bytes before relief. It remains one closed manifold genus-one component with no detected self-intersections
or zero-area triangles. Final volume is 6710837.897365 mm$^3$, removing 0.604764\% steel relative to the
previous closed mesh. Volume change against the original annularly repaired reference is -0.620987\%. The
original engineering surface remains open and has no valid mass. Maximum sampled distance from that source
remains 0.745438 mm, within the existing 0.75 mm guard. The current builder reproduces the relieved STL byte
for byte.

Independent clipping of neighboring triangles to their shared axial intervals gives minimum clearance to all
48 socket-wall triangle planes of 0.100019 mm at the adapter and 0.100022 mm at Segment 2. Triangle
collisions and bidirectional vertex containment are zero for both pairs. All eleven source CAD
configurations containing the three volumes use the same identity placements, so the same joint checks apply
to each. The separate 10,000-point Geant4 assembly test gives eight overlap reports before relief and no
exceptions afterward. Full default beamline construction now finishes without missing meshes, solid
warnings, or overlap reports. These checks cover the tested configuration and do not certify every possible
complete detector assembly.

\begin{table}[htbp]
\centering
\caption{Segment 1 before and after socket relief.}
\label{tab:downstream-joints-mesh}
\setlength{\tabcolsep}{9pt}
\begin{tabular}{lrrr}
\toprule
Quantity & Previous & Relieved & Change \\
\midrule
Triangles & 5,000 & 5,096 & +96 \\
Bytes & 250,084 & 254,884 & +4,800 \\
Boundary edges & 0 & 0 & 0 \\
Volume (mm$^3$) & 6751669.593 & 6710837.897 & -0.604764\% \\
\bottomrule
\end{tabular}
\end{table}

A new tracking comparison includes the adapter and both downstream segments in air. Its baseline is the
previous optimized closed meshes, not the original engineering inputs. The fixed GEMC3 executable,
materials, identity placements, matched 5 GeV target-origin primaries over 0.1--1.3 degrees and full
azimuth, FTFP\_BERT, no optics or field, one thread, default production cuts, and disabled GUI/event output
are identical. Five alternating-order repetitions use seeds 12345--12349 and 200 events per electron,
proton, positive-pion, and neutron run. All 52 processes finish normally and all forty transport runs report
200 events.

Previous/relieved median transport ratios are 0.981, 1.000, 1.008, 0.960 for electrons, protons, positive
pions, and neutrons respectively. Ratios above one denote faster tracking. The previous assembly produces 18
GeomNav1002, 4 GeomVol1002 exception blocks across the recorded processes; the relieved assembly produces 15
GeomNav1002 blocks. The remaining navigation warnings occur in the unchanged bellows region rather than
either relieved joint. Their z coordinates lie between 10177.906421 and 10281.864247 mm; all 1,741 triangles
intersecting the combined before/after warning interval are identical. Clean overlap checks therefore do not
establish warning-free tracking. These timings include overlap removal, local steel removal, and 96
additional facets. They measure the repair workload rather than a pure decimation gain, and do not establish
full CLAS12 performance or physics-response equivalence. The earlier study chapters and archives retain
their historical measurements.

\begin{table}[htbp]
\centering
\caption{Median transport wall seconds per 200 events: three-part downstream assembly.}
\label{tab:downstream-joints-tracking}
\setlength{\tabcolsep}{9pt}
\begin{tabular}{lrrr}
\toprule
Particle & Previous & Relieved & Previous/relieved \\
\midrule
e- & 11.25010 & 11.46960 & 0.981 \\
proton & 1.75729 & 1.75794 & 1.000 \\
pi+ & 1.93580 & 1.92101 & 1.008 \\
neutron & 1.09550 & 1.14060 & 0.960 \\
\bottomrule
\end{tabular}
\end{table}

\FloatBarrier
\begin{figure}[htbp]
\centering
\includegraphics[width=0.95\textwidth]{figures/downstream_joint_clearance.pdf}
\caption{Actual socket and neighboring tube polygons projected onto the most restrictive socket face at each
joint. The previous socket cuts into its neighbor; the relieved socket provides at least 0.1 mm face
clearance while the neighboring tube remains unchanged.}
\label{fig:downstream-joints}
\end{figure}

\FloatBarrier
\section{Beamline database run selection}
\label{sec:beamline-run-selection}

The beamline SQLite generator previously stored source labels such as \texttt{rgc\_summer2022} in the
variation column. The steering card selected \texttt{default}, so requesting run 16000 matched no geometry
rows and constructed only the world. The corrected generator publishes all beamline geometry and material
banks under \texttt{default}, indexed by run number. Source labels remain authoring selectors and ASCII
reference filenames; they are not additional SQLite variations.

With the regenerated local database, the command
\texttt{gemc beamline.yaml -gui -runno=16000} selects the summer-2022 configuration: ten native volumes
and 22 CAD volumes. A matched headless construction changes from one Geant4 volume before the fix to 33
afterward, including the world. The CLI and steering card need no variation override. Existing databases
require regeneration.

The three transverse CAD choices now use distinct run keys: 10000204 for \texttt{noWedges}, 10000206 for
\texttt{forwardWedges}, and 10000207 for \texttt{midpointWedges}. Each uses \texttt{default} and has its own
copies of the same two native pipe volumes and their materials. The dedicated transverse card selects the
companion CAD system. Regeneration replaces only the beamline systems in the selected experiment and removes
stale named banks; unrelated detector rows in an aggregate database are retained.

All 17 configured beamline runs construct with the expected SQLite row counts. The semantic comparator
passes all 15 native-geometry, 15 material, and 14 CAD comparisons against the GEMC2 source. A direct
database comparison preserves all physical fields, allowing only the transverse run-key changes and
the additional native copies. Repeated generation produces identical rows apart from database IDs,
and an aggregate-database check preserves unrelated rows exactly.

This is a database-selection repair. STL files, placements, material definitions, and the measured mesh
optimizations are unchanged; no new tracking-speed measurement is attributed to it. The retained
\texttt{data/beamline\_run\_selection/} evidence includes before/after database fixtures, the original
command's construction logs, all-run construction and comparator logs, validation results, and provenance
hashes. Historical benchmark archives retain their original inputs and source labels.

\FloatBarrier
\section{Moller cone tip}
\label{sec:cone-tip}

The eleventh beamline mesh study renames \texttt{cone-tip} to \texttt{moller\_cone\_tip} in all four source,
runtime, and GEMC3 CAD configurations: runs 5874, 11620, 18304, and 20000. The \texttt{beamline\_W} material,
passive 180-degree y rotation, and 1273.27 mm z translation remain unchanged. SQLite banks continue to use
\texttt{default}, indexed by run number. The original is archived; the three installed meshes have identical
bytes.

The original is a clean closed manifold solid with one component and genus two. Protected quadric decimation
reduces 13,730 triangles to 5,000 (63.5834\% fewer), and 686,584 bytes to 250,084 (63.5756\% fewer). All 1,936
faces adjoining edges above five degrees retain their exact triangle coordinates. The original bounds, central
passage, and other opening are retained; no hole filling or reconstruction is applied. The result has no
boundary edges, non-manifold elements, detected self-intersections, or zero-area triangles. The current
builder reproduces the saved STL byte for byte.

Volume changes from 196511.847100 to 196510.892015 mm$^3$, a change of -0.000486019\%. Bidirectional vertex,
edge, and surface sampling gives a maximum deviation of 0.032396 mm, below the 0.05 mm guard. This is a
sampled maximum, not a proven global error bound. More aggressive protected trials distort the geometry and
are rejected. Table~\ref{tab:cone-tip-mesh} gives the retained mesh properties.

\begin{table}[htbp]
\centering
\caption{Moller cone tip before and after protected decimation.}
\label{tab:cone-tip-mesh}
\begin{tabular}{lrr}
\toprule
Quantity & Original & Optimized \\
\midrule
Triangles & 13,730 & 5,000 \\
Bytes & 686,584 & 250,084 \\
Boundary edges & 0 & 0 \\
Components / genus & 1 / 2 & 1 / 2 \\
Volume (mm$^3$) & 196511.847100 & 196510.892015 \\
Detected intersecting faces & 0 & 0 \\
\bottomrule
\end{tabular}
\end{table}

The isolated benchmark selects the run-5874 CAD row with \texttt{default}; the tip is absent from run 11. The
benchmark runner now accepts \texttt{--runno} so it can select such components without remapping their run
keys. Both fixtures use matched 5 GeV electron, proton, positive-pion, and neutron primaries from the target,
uniform in solid angle over 2.5--4.8 degrees and full azimuth. Physics is FTFP\_BERT with no magnetic field or
optical processes, one thread, default production cuts, and no GUI or event output. Five alternating-order
repetitions use 200 events per particle and seeds 12345--12349. The fixed GEMC3 executable and its SHA-256
match the recent beamline studies; all primary hashes match within seed pairs.

Table~\ref{tab:cone-tip-tracking} reports median Geant4 beamOn wall times, excluding initialization. Ratios
are original/optimized; values above one mean faster tracking. All five pairs are faster for each particle,
and all 52 isolated benchmark processes complete without Geant4 exceptions. The sample measures this one mesh;
it does not establish complete CLAS12 speedup or detector-response equivalence.

\begin{table}[htbp]
\centering
\caption{Isolated Moller cone tip tracking; five repetitions of 200 events per particle.}
\label{tab:cone-tip-tracking}
\setlength{\tabcolsep}{6pt}
\begin{tabular}{lrrrr}
\toprule
Particle & Original (s) & Optimized (s) & Speedup & Paired ratio range \\
\midrule
Electron & 17.63610 & 10.63150 & 1.659 & 1.658--1.665 \\
Proton & 1.84329 & 1.62666 & 1.133 & 1.133--1.211 \\
Positive pion & 2.40739 & 2.04576 & 1.177 & 1.173--1.236 \\
Neutron & 1.62514 & 1.40266 & 1.159 & 1.159--1.207 \\
\bottomrule
\end{tabular}
\end{table}

Median loading process wall time is nearly unchanged, 0.843174 to 0.841411 s, while loading CPU time falls
from 0.80 to 0.76 s. Engine initialization wall time changes from 0.758880 to 0.727224 s. Startup peak RSS
falls from 133.125 to 122.813 MiB (7.75\%). Process CPU time, total wall time, memory, seed-paired ratios, and
raw diagnostics are retained separately from the tracking measurements.

The main-cone/tip assembly already overlaps before decimation. Independent surface collisions are present with
both engineering originals as well as with the current main cone. A separate 10,000-point Geant4 assembly
check reports five GeomVol1002 exception blocks both before and after the tip change; sampled penetrations
reach about 0.372 mm. This optimization preserves the engineering interfaces and does not repair that assembly
overlap. Collision counts and vertex-containment counts depend on tessellation and must not be interpreted as
overlap volumes or as evidence of improved clearance.

All four complete affected configurations construct 39 Geant4 volumes including the world. Their two
GeomSolids1001 warnings remain, and sampled overlap blocks change from 20 to 21 per construction. The renamed
volume changes construction order and the resulting sampled reports; contacts among unchanged parts remain.
The clean isolated benchmark therefore does not certify the complete assemblies. All 17 beamline
configurations pass volume-count construction checks. All 15 native-geometry, 15 material, and 14 CAD semantic
comparisons pass; direct database checks allow only the four names and STL paths to change. Every other mesh
retains its prior hash. Figure~\ref{fig:cone-tip} uses matched camera and bounds.

\begin{figure}[p]
\centering
\includegraphics[width=\linewidth]{figures/moller_cone_tip_stl_comparison.pdf}
\caption{Original and optimized Moller cone tip meshes with identical bounds and camera. Openings and
sharp features are retained while redundant smooth-surface triangles are removed.}
\label{fig:cone-tip}
\end{figure}

\FloatBarrier
\section{FT flange}
\label{sec:ft-flange}

The twelfth beamline mesh study renames \texttt{flange-assembly} to \texttt{FTFlange} in four configurations:
runs 5874, 11620, 18304, and 20000. The stainless-steel material, passive 180-degree y rotation, and 1273.27
mm z translation remain unchanged. Both repositories and the GEMC2 CAD database use the new name. GEMC3 SQLite
banks remain in \texttt{default}, selected by run number. The original is archived, and all three installed
meshes are byte-identical.

The original is already a clean closed annulus with only 552 triangles, one component, and genus one.
Topology-preserving quadric decimation reduces it to 540 triangles (2.1739\% fewer); stored bytes fall from
27,684 to 27,084 (2.1673\% fewer). No holes are filled or surfaces reconstructed. The central passage and
overall bounds are retained. After edge collapse, vertices are restored to the two exact engineering face
planes, local z = -526.342956543 and -519.992980957 mm. The resulting STL has no boundary edges, non-manifold
elements, detected self-intersections, or zero-area triangles, and the builder reproduces it byte for byte.

Steel volume changes from 229348.479344 to 229348.228841 mm$^3$ (-0.000109224\%). Bidirectional vertex, edge,
and surface sampling gives a maximum deviation of 0.049822 mm, below the 0.05 mm guard. This sampled maximum
is not a proven global error bound. A 536-triangle trial exceeds that guard; larger reductions are rejected.
The small reduction reflects the already-coarse engineering tessellation. Table~\ref{tab:ft-flange-mesh}
summarizes the mesh changes, and Figure~\ref{fig:ft-flange} compares the stored triangles with a matched
camera.

The isolated benchmark selects run 5874 with \texttt{default}, using matched 5 GeV electron, proton,
positive-pion, and neutron primaries from the target, uniform in solid angle over 1.4--3.8 degrees and full
azimuth. Physics is FTFP\_BERT with no field or optical processes, one thread, default production cuts, and no
GUI or event output. Five alternating-order repetitions use 2,000 events per particle and seeds 12345--12349.
The executable matches the recent beamline studies, and all primary hashes agree within seed pairs.
Table~\ref{tab:ft-flange-tracking} reports median Geant4 beamOn wall times excluding initialization, with
original/optimized ratios and the ranges of paired ratios.

Tracking is effectively unchanged: median ratios are 1.0018 for electrons, 1.0001 for protons, 1.0003 for
positive pions, and 0.9993 for neutrons. Every paired range includes one; only three or four of the five pairs
are faster. These variations do not establish a tracking-speed gain. All 52 isolated processes complete
without Geant4 exceptions. Median loading process wall time changes from 0.784696 to 0.791953 s, CPU time
remains 0.72 s, and engine initialization wall time changes from 0.681824 to 0.672900 s. Startup peak RSS is
115.391 versus 115.375 MiB, effectively unchanged. Complete process CPU, wall time, memory, and per-seed
results are retained separately from the transport measurements.

The neighboring inner and outer assemblies already contact or overlap the flange before this change.
Independent triangle collision counts are 449 versus 449 for the inner assembly and 600 versus 583 for the
outer assembly; those tessellation-dependent counts are not overlap volumes or clearance proofs. The unchanged
inner assembly is not closed. A separate three-part 10,000-point Geant4 check reports one GeomSolids1001 and
eight GeomVol1002 exception blocks both before and after optimization. Sampled penetration depths reach about
0.249 mm. This mesh simplification retains the existing engineering interfaces; the assembly defects remain.

All four complete affected configurations construct 39 Geant4 volumes including the world. Both versions
produce two GeomSolids1001 and 21 GeomVol1002 blocks per complete construction. All 17 beamline configurations
pass volume-count construction checks, and the comparator passes 15 native geometry, 15 material, and 14 CAD
comparisons. Direct database checks allow only the four names and STL paths to change; every other mesh
retains its previous hash. The clean isolated benchmark does not certify complete geometry or
detector-response equivalence.

\begin{table}[htbp]
\centering
\caption{FT flange before and after conservative decimation.}
\label{tab:ft-flange-mesh}
\setlength{\tabcolsep}{6pt}
\begin{tabular}{lrr}
\toprule
Quantity & Original & Optimized \\
\midrule
Triangles & 552 & 540 \\
Bytes & 27,684 & 27,084 \\
Boundary edges & 0 & 0 \\
Components / genus & 1 / 1 & 1 / 1 \\
Volume (mm$^3$) & 229348.479344 & 229348.228841 \\
\bottomrule
\end{tabular}
\end{table}

\begin{table}[htbp]
\centering
\caption{Isolated FT flange tracking; five repetitions of 2,000 events per particle.}
\label{tab:ft-flange-tracking}
\setlength{\tabcolsep}{6pt}
\begin{tabular}{lrrrr}
\toprule
Particle & Original (s) & Optimized (s) & Ratio & Paired ratio range \\
\midrule
Electron & 2.68847 & 2.68375 & 1.0018 & 0.9998--1.0054 \\
Proton & 3.61487 & 3.61435 & 1.0001 & 0.9957--1.0029 \\
Positive pion & 3.06986 & 3.06909 & 1.0003 & 0.9994--1.0033 \\
Neutron & 1.75223 & 1.75338 & 0.9993 & 0.9986--1.0051 \\
\bottomrule
\end{tabular}
\end{table}

\begin{figure}[p]
\centering
\includegraphics[width=\linewidth]{figures/ft_flange_stl_comparison.pdf}
\caption{Original and optimized FT flange meshes with identical bounds and camera. The central passage, flat
axial faces, and overall shape are retained; the already-coarse tessellation changes modestly.}
\label{fig:ft-flange}
\end{figure}

\FloatBarrier
\section{Inner Moller-cone extension}
\label{sec:extension-inner}

The thirteenth beamline mesh study renames \texttt{inner-assembly} to
\texttt{moller-\allowbreak{}cone-\allowbreak{}extension-\allowbreak{}inner} in
runs 5874, 11620, 18304, and 20000. Both repositories, their CAD placements, and the GEMC2 CAD database use
the new name. Stainless steel, the passive 180-degree y rotation, and the 1273.27 mm z translation remain
unchanged. GEMC3 SQLite banks remain under \texttt{default}, indexed by run number. The engineering original
is archived, and the three installed optimized STLs have identical bytes.

The original has 18,057 triangles, 75 boundary edges in two loops, one connected component, and no detected
self-intersections or zero-area triangles. The two open rims surround one radial port on the positive local-x
side: a 35-vertex rim at tube radius 40 mm and a 40-vertex rim at radius 47.625 mm. They share a nominal 6.35
mm port radius centered at local z = -418.042413330 mm. The wall between these rims is absent. Stitching their
angularly ordered boundaries adds 75 triangles while retaining every original vertex. The beam passage and
both radial ports remain open; no hole caps or surface reconstruction are applied.

This repair produces an 18,132-triangle closed reference with one manifold component and genus three. Genus
describes the closed solid after restoring the missing tunnel wall; the open original does not define a valid
solid volume. Protected quadric decimation retains all 7,297 sharp-edge and port-adjacent faces at their exact
coordinates and reduces the result to 12,500 triangles. This is 30.7748\% fewer than the raw source and
31.0611\% fewer than the repaired reference. Original bounds are retained. The final mesh has no boundary
edges, non-manifold elements, detected self-intersections, or zero-area triangles. The minimum triangle
altitude is 0.002261 mm. Rebuilding reproduces the reference and optimized STL bytes exactly.
Table~\ref{tab:extension-inner-mesh} gives the mesh inventory, and Figure~\ref{fig:extension-inner} compares
the stored meshes.

Repaired reference volume is 1431693.961982 mm$^3$ and optimized volume is 1431694.095171 mm$^3$, a
+0.000009303\% decimation change. The original open mesh has no valid engineering mass. Maximum bidirectional
vertex/edge/surface sampling between the repaired reference and optimized STL is 0.015669 mm, below the 0.05
mm guard; this is a sampled measurement, not a global bound or a distance to the incomplete raw surface.
Larger reductions distort the geometry and are rejected. Independent containment probes check 25 points along
the beam passage and six points in and around the repaired port, confirming that the passage and port remain
empty and adjacent steel remains filled.

Two separate isolated benchmarks use run 5874 with \texttt{default}, matched 5 GeV electron, proton,
positive-pion, and neutron primaries from the target, uniform in solid angle over 0.8--1.8 degrees and full
azimuth. Physics is FTFP\_BERT with no field or optical processes, one thread, default production cuts, and no
GUI or event output. Each benchmark uses five alternating-order repetitions of 200 events per particle with
seeds 12345--12349. The fixed executable matches recent beamline studies; every seed pair has matching primary
hashes.

Table~\ref{tab:extension-inner-upgrade} compares the defective raw mesh with the upgraded mesh.
Table~\ref{tab:extension-inner-control} separately compares the closed, unreduced reference with the optimized
mesh, preserving the repaired shape and material volume to the reported decimation precision. Timings are
median Geant4 beamOn wall times excluding initialization; ratios are baseline/optimized, so values above one
indicate faster tracking. Raw timings include surface and navigation defects and do not establish faster
correct transport. Only the repaired-reference control measures tessellation changes after the missing wall is
restored.

Raw-upgrade tracking ratios are 1.138, 1.055, 1.088, 1.043 for electrons, protons, positive pions, and
neutrons, respectively. Loading process wall time changes from 0.937190 to 0.886479 s; loading CPU time is
0.85 versus 0.83 s. Startup peak RSS is 141.844 versus 134.578 MiB. Complete process CPU, memory, wall time,
and seed-paired results are retained independently.

Repaired-reference tracking ratios are 1.261, 1.086, 1.144, 1.095 for electrons, protons, positive pions, and
neutrons, respectively. Loading process wall time changes from 0.892270 to 0.841108 s; loading CPU time is
0.82 versus 0.78 s. Startup peak RSS is 141.031 versus 133.953 MiB. Complete process CPU, memory, wall time,
and seed-paired results are retained independently.

Across the valid raw-input benchmark processes, exception totals are 163 GeomNav1002, 26 GeomSolids1001. All
optimized runs and both sides of the repaired-reference control have no Geant4 exceptions. These are
process-level diagnostic counts, not unique physical overlaps or counts of affected events.

A separate six-part 10,000-point assembly check includes the extension, FT flange, FT inner shield, main cone,
lead shield, and outer assembly. Exception blocks are original: 1 GeomSolids1001, 15 GeomVol1002; reference:
15 GeomVol1002; optimized: 15 GeomVol1002. The restored wall removes the extension surface-hole warning;
engineering overlaps remain. Independent surface collisions use the actual CAD transforms, including the FT
inner shield offset. Collision counts depend on tessellation and are not overlap volumes or clearance proofs.

All four complete affected configurations construct 39 Geant4 volumes including the world. Per-configuration
exception blocks are original: 2 GeomSolids1001, 21 GeomVol1002; optimized: 1 GeomSolids1001, 20 GeomVol1002.
All 17 beamline configurations pass volume-count construction checks. All 15 native geometry, 15 material, and
14 CAD semantic comparisons pass. Direct database checks allow only the four names and STL paths to change;
every other mesh retains its prior hash. These checks do not establish complete overlap freedom, full CLAS12
speed gains, or detector-response equivalence.

\begin{table}[htbp]
\centering
\caption{Inner extension raw input, repaired reference, and optimized mesh.}
\label{tab:extension-inner-mesh}
\setlength{\tabcolsep}{5pt}
\begin{tabular}{lrrr}
\toprule
Quantity & Raw & Repaired & Optimized \\
\midrule
Triangles & 18,057 & 18,132 & 12,500 \\
Bytes & 902,934 & 906,684 & 625,084 \\
Boundary edges & 75 & 0 & 0 \\
Closed components / genus & undefined & 1 / 3 & 1 / 3 \\
Volume (mm$^3$) & undefined & 1431693.961982 & 1431694.095171 \\
\bottomrule
\end{tabular}
\end{table}

\begin{table}[htbp]
\centering
\caption{Raw versus optimized inner extension; five repetitions of 200 events per particle.}
\label{tab:extension-inner-upgrade}
\setlength{\tabcolsep}{5pt}
\begin{tabular}{lrrrr}
\toprule
Particle & Raw (s) & Optimized (s) & Ratio & Paired ratio range \\
\midrule
Electron & 26.50810 & 23.28480 & 1.138 & 1.117--1.145 \\
Proton & 2.90723 & 2.75559 & 1.055 & 1.034--1.058 \\
Positive pion & 4.34586 & 3.99316 & 1.088 & 1.065--1.113 \\
Neutron & 2.63253 & 2.52499 & 1.043 & 1.035--1.068 \\
\bottomrule
\end{tabular}
\end{table}

\begin{table}[htbp]
\centering
\caption{Repaired versus optimized inner extension; five repetitions of 200 events per particle.}
\label{tab:extension-inner-control}
\setlength{\tabcolsep}{5pt}
\begin{tabular}{lrrrr}
\toprule
Particle & Repaired (s) & Optimized (s) & Ratio & Paired ratio range \\
\midrule
Electron & 29.39720 & 23.31290 & 1.261 & 1.214--1.275 \\
Proton & 3.00190 & 2.76398 & 1.086 & 1.082--1.116 \\
Positive pion & 4.56824 & 3.99481 & 1.144 & 1.130--1.144 \\
Neutron & 2.76308 & 2.52386 & 1.095 & 1.091--1.118 \\
\bottomrule
\end{tabular}
\end{table}

\begin{figure}[p]
\centering
\includegraphics[width=\linewidth]{figures/extension_inner_stl_comparison.pdf}
\caption{Original and optimized inner Moller-cone extension meshes with matched bounds and camera. The beam
passage and both radial ports remain open; the missing wall of the lower port is restored.}
\label{fig:extension-inner}
\end{figure}

\FloatBarrier
\section{Moller-cone extension lead fill}
\label{sec:extension_lead}

This study renames \texttt{lead-\allowbreak{}shield} to
\texttt{moller-\allowbreak{}cone-\allowbreak{}extension-\allowbreak{}lead-\allowbreak{}fill} in runs 5874,
11620, 18304, and 20000 in both repositories. The lead material, passive 180-degree y rotation, and 1273.27 mm
z translation are retained. SQLite banks remain in \texttt{default}, selected by run number. The original is
archived, and all three installed optimized meshes are byte-identical.

The original is already a closed manifold solid with 10,594 triangles, one component, and genus two. Protected
quadric decimation reduces it to 8,000 triangles (24.4856\% fewer), with 529,784 versus 400,084 stored bytes.
All 7,156 faces adjoining sharp edges above five degrees or any vertex within radius 48 mm retain their exact
triangle coordinates. This protects the full inner mating surface and adjacent end faces. The beam passage,
radial opening, and exact bounds remain. There is no hole filling or surface reconstruction. The result
retains its topology and has no boundary edges, non-manifold elements, detected intersections, or zero-area
triangles. Independent rebuilding reproduces its bytes.

Material volume changes from 11950334.138104 to 11950334.137844 mm$^3$ (-0.0000000022\%). Maximum
bidirectional vertex, edge, and surface sampled deviation is 0.00001998 mm, below the 0.05 mm guard. This
sampled maximum is not a proven global bound. Trials exceeding the shape, bounds, or volume guards are
rejected.

The isolated benchmark selects run 5874 with \texttt{default}. It uses matched 5 GeV electrons, protons,
positive pions, and neutrons from the target, uniform in solid angle over 1.6--4.5 degrees and full azimuth.
Physics is FTFP\_BERT with no field or optical processes, one thread, default production cuts, and no GUI or
event output. Five alternating-order repetitions use 200 events per particle and seeds 12345--12349. The fixed
executable matches the recent beamline studies, and primary hashes agree within each seed pair.

Tracking tables report median Geant4 beamOn wall time excluding initialization. Ratios are original/optimized;
values above one indicate faster tracking. Median ratios are 1.0826, 1.0795, 1.1226, 1.0882 for electrons,
protons, positive pions, and neutrons, respectively; faster-pair counts are 5/5, 5/5, 5/5, 5/5. Every paired
measurement improves in this sample. All 52 isolated processes complete without Geant4 exceptions. These
measurements describe this isolated workload and do not establish complete CLAS12 performance or
detector-response equivalence.

Loading medians, in seconds, are process wall time: 0.896672 versus 0.844312; process CPU time: 0.850000
versus 0.810000; engine initialization wall time: 0.803922 versus 0.767714. Startup peak RSS is 137.297 versus
133.859 MiB. Process wall time includes launch and initialization; it is separate from beamOn transport.
Complete process CPU, wall time, peak memory, per-seed timings, and provenance are retained.

The separate four-part, 10,000-point assembly check uses the fill, casing, inner extension, and FT flange.
Existing assembly overlaps remain. Original and optimized exception counts are 10 GeomVol1002 and 10
GeomVol1002. Independent triangle-contact counts are moller-cone-extension-inner: 10273 versus 10273;
outer-assembly: 15473 versus 11052; FTFlange: 819 versus 819. These counts depend on tessellation and are
neither penetration volumes nor clearance proofs. Protected mating faces prevent this optimization from moving
those interfaces; retaining an interface does not repair pre-existing interference.

All four complete affected configurations construct 39 Geant4 volumes including the world. Run 5874
exception-block counts are original: 1 GeomSolids1001, 20 GeomVol1002; optimized: 1 GeomSolids1001, 21
GeomVol1002. The archive reports all four runs individually. All 17 beamline configurations pass volume-count
construction checks, and semantic comparisons pass 15 native geometry, 15 material, and 14 CAD comparisons.
Direct database checks allow only four names and STL paths to change; every other mesh retains its previous
hash. Existing complete-assembly warnings are separate from the clean isolated timings. A separate rename-only
control retains the original fill STL under the new name and path and produces the same 1 GeomSolids1001 and
21 GeomVol1002 blocks as the optimized configuration. The count increase from 20 to 21 also occurs without
decimation; overlap-check ordering and sampling change after the rename.

\begin{table}[htbp]
\centering
\caption{Moller-cone extension lead fill mesh inventory.}
\label{tab:extension_lead0}
\setlength{\tabcolsep}{5pt}
\begin{tabular}{lrr}
\toprule
Quantity & Original & Optimized \\
\midrule
Triangles & 10,594 & 8,000 \\
Bytes & 529,784 & 400,084 \\
Boundary edges & 0 & 0 \\
Components / genus & 1 / 2 & 1 / 2 \\
Volume (mm$^3$) & 11950334.138104 & 11950334.137844 \\
\bottomrule
\end{tabular}
\end{table}

\begin{table}[htbp]
\centering
\caption{Moller-cone extension lead fill tracking; five repetitions of 200 events per particle.}
\label{tab:extension_lead1}
\setlength{\tabcolsep}{5pt}
\begin{tabular}{lrrrr}
\toprule
Particle & Before (s) & After (s) & Ratio & Paired range \\
\midrule
Electron & 23.76520 & 21.95230 & 1.0826 & 1.0810--1.1063 \\
Proton & 4.70122 & 4.35507 & 1.0795 & 1.0672--1.0912 \\
Positive pion & 7.60341 & 6.77309 & 1.1226 & 1.0457--1.1318 \\
Neutron & 5.23546 & 4.81107 & 1.0882 & 1.0729--1.1069 \\
\bottomrule
\end{tabular}
\end{table}

\begin{figure}[p]
\centering
\includegraphics[width=\linewidth]{figures/extension_lead_stl_comparison.pdf}
\caption{Moller-cone extension lead fill before and after optimization, with identical bounds and camera. The
bore and
radial opening remain.}
\label{fig:extension_lead}
\end{figure}

\FloatBarrier
\section{Moller-cone extension lead case}
\label{sec:extension_case}

This study renames \texttt{outer-\allowbreak{}assembly} to
\texttt{moller-\allowbreak{}cone-\allowbreak{}extension-\allowbreak{}lead-\allowbreak{}case} in runs 5874,
11620, 18304, and 20000 in both repositories. The stainless steel material, passive 180-degree y rotation, and
1273.27 mm z translation are retained. SQLite banks remain in \texttt{default}, selected by run number. The
original is archived, and all three installed optimized meshes are byte-identical.

The original is already a closed manifold solid with 17,360 triangles, one component, and genus two. Protected
quadric decimation reduces it to 10,000 triangles (42.3963\% fewer), with 868,084 versus 500,084 stored bytes.
All 9,105 faces adjoining sharp edges above five degrees or any vertex within radius 118 mm retain their exact
triangle coordinates. This protects the full inner mating surface and adjacent end faces. The beam passage,
radial opening, and exact bounds remain. There is no hole filling or surface reconstruction. The result
retains its topology and has no boundary edges, non-manifold elements, detected intersections, or zero-area
triangles. Independent rebuilding reproduces its bytes.

Material volume changes from 817759.770060 to 817759.420845 mm$^3$ (-0.0000427038\%). Maximum bidirectional
vertex, edge, and surface sampled deviation is 0.00678680 mm, below the 0.05 mm guard. This sampled maximum is
not a proven global bound. Trials exceeding the shape, bounds, or volume guards are rejected.

The isolated benchmark selects run 5874 with \texttt{default}. It uses matched 5 GeV electrons, protons,
positive pions, and neutrons from the target, uniform in solid angle over 3.7--4.8 degrees and full azimuth.
Physics is FTFP\_BERT with no field or optical processes, one thread, default production cuts, and no GUI or
event output. Five alternating-order repetitions use 200 events per particle and seeds 12345--12349. The fixed
executable matches the recent beamline studies, and primary hashes agree within each seed pair.

Tracking tables report median Geant4 beamOn wall time excluding initialization. Ratios are original/optimized;
values above one indicate faster tracking. Median ratios are 1.0529, 1.0377, 1.0098, 1.0198 for electrons,
protons, positive pions, and neutrons, respectively; faster-pair counts are 5/5, 5/5, 5/5, 2/5. Some paired
ranges include one; those particle samples do not establish a consistent gain. All 52 isolated processes
complete without Geant4 exceptions. These measurements describe this isolated workload and do not establish
complete CLAS12 performance or detector-response equivalence.

Loading medians, in seconds, are process wall time: 1.000175 versus 0.838249; process CPU time: 0.920000
versus 0.760000; engine initialization wall time: 0.882605 versus 0.724002. Startup peak RSS is 137.469 versus
127.516 MiB. Process wall time includes launch and initialization; it is separate from beamOn transport.
Complete process CPU, wall time, peak memory, per-seed timings, and provenance are retained.

The separate four-part, 10,000-point assembly check uses the fill, casing, inner extension, and FT flange.
Existing assembly overlaps remain. Original and optimized exception counts are 9 GeomVol1002 and 9
GeomVol1002. Independent triangle-contact counts are moller-cone-extension-inner: 0 versus 0;
moller-cone-extension-lead-fill: 10990 versus 10990; FTFlange: 583 versus 583. These counts depend on
tessellation and are neither penetration volumes nor clearance proofs. Protected mating faces prevent this
optimization from moving those interfaces; retaining an interface does not repair pre-existing interference.

All four complete affected configurations construct 39 Geant4 volumes including the world. Run 5874
exception-block counts are original: 1 GeomSolids1001, 21 GeomVol1002; optimized: 1 GeomSolids1001, 21
GeomVol1002. The archive reports all four runs individually. All 17 beamline configurations pass volume-count
construction checks, and semantic comparisons pass 15 native geometry, 15 material, and 14 CAD comparisons.
Direct database checks allow only four names and STL paths to change; every other mesh retains its previous
hash. Existing complete-assembly warnings are separate from the clean isolated timings.

\begin{table}[htbp]
\centering
\caption{Moller-cone extension lead case mesh inventory.}
\label{tab:extension_case0}
\setlength{\tabcolsep}{5pt}
\begin{tabular}{lrr}
\toprule
Quantity & Original & Optimized \\
\midrule
Triangles & 17,360 & 10,000 \\
Bytes & 868,084 & 500,084 \\
Boundary edges & 0 & 0 \\
Components / genus & 1 / 2 & 1 / 2 \\
Volume (mm$^3$) & 817759.770060 & 817759.420845 \\
\bottomrule
\end{tabular}
\end{table}

\begin{table}[htbp]
\centering
\caption{Moller-cone extension lead case tracking; five repetitions of 200 events per particle.}
\label{tab:extension_case1}
\setlength{\tabcolsep}{5pt}
\begin{tabular}{lrrrr}
\toprule
Particle & Before (s) & After (s) & Ratio & Paired range \\
\midrule
Electron & 4.14328 & 3.93523 & 1.0529 & 1.0292--1.1052 \\
Proton & 0.99817 & 0.96193 & 1.0377 & 1.0001--1.0377 \\
Positive pion & 0.88047 & 0.87188 & 1.0098 & 1.0032--1.0308 \\
Neutron & 0.81340 & 0.79759 & 1.0198 & 0.9544--1.0198 \\
\bottomrule
\end{tabular}
\end{table}

\begin{figure}[p]
\centering
\includegraphics[width=\linewidth]{figures/extension_case_stl_comparison.pdf}
\caption{Moller-cone extension lead case before and after optimization, with identical bounds and camera. The
bore and
radial opening remain.}
\label{fig:extension_case}
\end{figure}

\FloatBarrier
\section{Moller-cone extension tungsten add-on}
\label{sec:extension-tungsten}

The sixteenth beamline mesh study renames \texttt{outer-\allowbreak{}cone} to
\texttt{moller-\allowbreak{}cone-\allowbreak{}extension-\allowbreak{}tungsten-\allowbreak{}addon} in runs
5874, 11620, 18304, and 20000. Both repositories, their CAD placements, and the GEMC2 CAD database use the new
name. The \texttt{beamline\_W} material, passive 180-degree y rotation, and 1273.27 mm z translation remain
unchanged. GEMC3 banks remain under \texttt{default}, indexed by run number. The original is archived, and all
three optimized meshes have identical bytes.

The original has 27,778 triangles and 68 boundary edges in two loops. These 33- and 35-vertex rims surround an
inclined radial port whose connecting wall is missing. Fitting the rim coordinates indicates a 6.350052 mm
radius and a 4.499946-degree inclination toward positive local z along positive local y. The builder orders
the existing rim vertices around the nominal 4.5-degree axis; it imposes no new rim coordinates. Bridging them
adds 68 wall triangles without moving any original vertex or removing any original face geometry. No beam
passage or radial opening is capped; there is no surface reconstruction. The independent repair audit retains
the original vertices and triangle coordinates exactly, allowing coherent face orientation.

The repaired reference is a closed manifold component with genus five and 27,846 triangles. The open source
has no valid solid volume; its reported genus four describes a surface with boundary, not the repaired solid.
Protected quadric decimation retains 9,154 faces, including sharp features, port-adjacent faces, faces with a
vertex within 0.25 mm of the main Moller cone, and faces adjoining the extension joint below local z = -169.5
mm. The main-cone dependency is guarded by its mesh hash. The final mesh has 12,000 triangles, 56.80\% fewer
than the raw source and 56.91\% fewer than the repaired reference. Its 600,084 bytes compare with 1,388,984
raw and 1,392,384 repaired bytes. Bounds and all protected face coordinates are exact. The mesh is closed,
manifold, and free of detected intersections and zero-area triangles. Table~\ref{tab:extension-tungsten-mesh}
inventories the meshes; Figure~\ref{fig:extension-tungsten} compares their tessellation.

Repaired tungsten volume is 3224441.660431 mm$^3$ and optimized volume is 3224441.511484 mm$^3$, a
-0.0000046193\% decimation change. Maximum bidirectional vertex/edge/surface sampled deviation from the
repaired reference is 0.012837 mm, below the 0.05 mm guard. This is a sampled measurement, not a proven global
bound. Minimum triangle altitude is 0.000473 mm. More aggressive protected trials distort the shape and are
rejected. Independent probes check 25 points along the beam passage, three along the inclined port, and four
in adjacent tungsten. The passages remain empty and the surrounding material remains filled. Both the repaired
and optimized STLs reproduce byte for byte.

Two separate isolated benchmarks select run 5874 with \texttt{default}. Matched 5 GeV electron, proton,
positive-pion, and neutron primaries originate at the target, uniform in solid angle over 2.7--5.0 degrees and
full azimuth. Physics is FTFP\_BERT with no field or optical processes, one thread, default production cuts,
and no GUI or event output. Each comparison uses five alternating-order repetitions of 200 events per
particle, seeds 12345--12349, the same fixed executable as the recent beamline studies, and identical primary
hashes within seed pairs.

Table~\ref{tab:extension-tungsten-upgrade} compares raw input with the upgraded mesh.
Table~\ref{tab:extension-tungsten-control} compares the closed, unreduced reference with the optimized mesh.
Median Geant4 beamOn wall times exclude initialization; ratios are baseline/optimized, so values above one
indicate faster tracking. Raw timings include the incomplete port wall and navigation defects and do not
establish faster correct transport. The repaired-reference control measures decimation with the repaired shape
and material volume preserved to the reported precision.

Raw-upgrade tracking ratios are 1.754, 1.488, 1.634, 1.472 for electrons, protons, positive pions, and
neutrons, respectively. Faster-pair counts are 5/5, 5/5, 5/5, 5/5 in the same order. Loading process wall time
changes from 0.997174 to 0.890545 s; loading CPU time is 0.93 versus 0.81 s. Startup peak RSS is 156.969
versus 131.500 MiB. Complete process CPU, memory, wall time, and seed-paired results are retained
independently.

Repaired-reference tracking ratios are 1.647, 1.548, 1.642, 1.535 for electrons, protons, positive pions, and
neutrons, respectively. Faster-pair counts are 5/5, 5/5, 5/5, 5/5 in the same order. Loading process wall time
changes from 0.951752 to 0.840506 s; loading CPU time is 0.90 versus 0.79 s. Startup peak RSS is 155.031
versus 131.188 MiB. Complete process CPU, memory, wall time, and seed-paired results are retained
independently.

Across the valid raw-input benchmark processes, exception totals are 66 GeomNav1002, 26 GeomSolids1001. All
optimized runs and both sides of the repaired-reference control have no Geant4 exceptions. These are
process-level diagnostic counts, not unique physical overlaps or counts of affected events.

A separate five-part 10,000-point assembly check includes the add-on, main cone, inner extension, lead fill,
and lead case. Exception blocks are original: 1 GeomSolids1001, 12 GeomVol1002; reference: 12 GeomVol1002;
optimized: 12 GeomVol1002. The restored wall removes the extension surface-hole warning; engineering overlaps
remain. Independent surface collisions use the actual CAD transforms, with identical passive placements.
Main-cone contact counts remain 3041 and lead-case counts remain 233 in all three fixtures; the inner
extension and lead fill have zero surface contacts with the add-on. Collision counts depend on tessellation
and are not overlap volumes or clearance proofs.

All four complete affected configurations construct 39 Geant4 volumes including the world. Per-configuration
exception blocks are original: 1 GeomSolids1001, 21 GeomVol1002; optimized: 21 GeomVol1002. All 17 beamline
configurations pass volume-count construction checks. All 15 native geometry, 15 material, and 14 CAD semantic
comparisons pass. Direct database checks allow only the four names and STL paths to change; every other mesh
retains its prior hash. These checks do not establish complete overlap freedom, full CLAS12 speed gains, or
detector-response equivalence.

\begin{table}[htbp]
\centering
\caption{Tungsten add-on raw input, repaired reference, and optimized mesh.}
\label{tab:extension-tungsten-mesh}
\setlength{\tabcolsep}{5pt}
\begin{tabular}{lrrr}
\toprule
Quantity & Raw & Repaired & Optimized \\
\midrule
Triangles & 27,778 & 27,846 & 12,000 \\
Bytes & 1,388,984 & 1,392,384 & 600,084 \\
Boundary edges & 68 & 0 & 0 \\
Closed components / genus & undefined & 1 / 5 & 1 / 5 \\
Volume (mm$^3$) & undefined & 3224441.660431 & 3224441.511484 \\
\bottomrule
\end{tabular}
\end{table}

\begin{table}[htbp]
\centering
\caption{Raw versus optimized tungsten add-on; five repetitions of 200 events per particle.}
\label{tab:extension-tungsten-upgrade}
\setlength{\tabcolsep}{5pt}
\begin{tabular}{lrrrr}
\toprule
Particle & Raw (s) & After (s) & Ratio & Paired range \\
\midrule
Electron & 26.35900 & 15.02950 & 1.754 & 1.735--1.781 \\
Proton & 4.75147 & 3.19228 & 1.488 & 1.409--1.491 \\
Positive pion & 7.16772 & 4.38598 & 1.634 & 1.523--1.634 \\
Neutron & 4.33537 & 2.94566 & 1.472 & 1.472--1.508 \\
\bottomrule
\end{tabular}
\end{table}

\begin{table}[htbp]
\centering
\caption{Repaired versus optimized tungsten add-on; five repetitions of 200 events per particle.}
\label{tab:extension-tungsten-control}
\setlength{\tabcolsep}{5pt}
\begin{tabular}{lrrrr}
\toprule
Particle & Repaired (s) & After (s) & Ratio & Paired range \\
\midrule
Electron & 27.83340 & 16.90260 & 1.647 & 1.554--1.874 \\
Proton & 5.06416 & 3.27214 & 1.548 & 1.496--1.564 \\
Positive pion & 7.66426 & 4.66753 & 1.642 & 1.568--1.642 \\
Neutron & 4.69515 & 3.05862 & 1.535 & 1.499--1.556 \\
\bottomrule
\end{tabular}
\end{table}

\begin{figure}[p]
\centering
\includegraphics[width=\linewidth]{figures/extension_tungsten_stl_comparison.pdf}
\caption{Original and optimized Moller-cone extension tungsten add-on meshes with matched bounds and camera.
The beam passage and radial ports remain open; the missing inclined port wall is restored.}
\label{fig:extension-tungsten}
\end{figure}

\FloatBarrier
\section{FT and Moller-cone interface clearances}
\label{sec:ft-joints}

This October 16 revision corrects the remaining FT/Moller-cone interfaces after the individual mesh studies.
Five meshes receive local material-removing operations: the inner extension, lead fill, lead case, main cone,
and tungsten add-on. FTInnerShield, FTFlange, and the cone tip remain byte-identical. Placements, passive
rotations, materials, names, and run keys are unchanged. All SQLite banks remain in \texttt{default}, indexed
by run number.

Several original engineering bores are coarse inscribed polygons. Their facets intrude into neighboring
surfaces even when nominal vertex radii agree. The inner extension also has a small shoulder at the main-cone
joint. The corrections use closed, polygonal cutters whose faces enclose the desired circular clearance
profiles, together with a conservative convex envelope of the actual lead fill. Cutters are intermediate
construction objects and are not additional detector volumes.

The inner-extension bore clears the placed FT shield, which is offset -369.73 mm in the common cone frame. Its
end is relieved around the unchanged flange bore. The lead-fill bore clears the inner extension. A scaled
convex envelope of the pre-clearance lead fill supplies at least 0.02002 mm plane clearance before final
casing simplification; a separate local cutter clears the unchanged flange. This envelope reduces casing
material loss from 4.4702\% in a rejected circular-envelope trial to 1.2093\%. The main cone receives a short
stepped relief around the inner-extension shoulder and a profiled opening around the cone tip. Its upper face
moves from z=712.002991 to 711.950012 mm. The add-on bore clears the main cone by a nominal 0.10 mm radial
profile margin. The health report retains every cutter profile and input hash.

All five corrected meshes retain one closed manifold component and their previous genera, with no detected
self-intersections or zero-area facets. The existing beam passages and radial ports remain open. A 0.0001 mm
weld of the saved float32 meshes removes numerical fragments. The casing is simplified from its boolean
reference to 16000 triangles with endpoint collapses and a maximum sampled deviation of 0.007648 mm; its
volume changes by +0.00002023\% during this simplification. It contains more triangles than the earlier
10000-triangle casing but remains below the engineering source. The eight-part total changes from 59540 to
60100 triangles (+0.9405\%). The material changes in Table~\ref{tab:ft-joints-meshes} are intentional
clearance removals, separate from the earlier decimation volume guards.

An initial cut end face passed overlap checks but produced nine GeomNav1002 boundary-navigation warnings
during tracking. Retriangulating that annular face as an angular strip retains its exact surface, boundary
coordinates, triangle count, and material volume. The previously failing electron seed passes after this
change. The initial benchmark is retained as an excluded trial and is not used in the final timing table.
Canonical STL output changes facet order and cyclic starting vertices without changing geometry or winding;
all eight outputs and the casing reference reproduce byte for byte.

Independent placed-surface checks find zero triangle contacts and zero contained vertices in both directions
at all eight corrected interfaces. Table~\ref{tab:ft-joints-gaps} gives minimum vertex-to-surface samples, not
proven global distance bounds. Three seeds (12345--12347) with 20000-point Geant4 overlap checks produce no
exceptions. The previous eight-part fixture produced 21 GeomVol1002 blocks. Matched meridian sections show the
local changes in Figures~\ref{fig:ft-joints-1} and~\ref{fig:ft-joints-2}; the FT shield uses its actual
placement offset.

The isolated tracking comparison includes the corrected eight-part assembly in run 5874 with \texttt{default}.
Matched 5 GeV electrons, protons, positive pions, and neutrons originate at the target, uniform in solid angle
over 0.8--4.8 degrees and full azimuth. FTFP\_BERT uses one thread, default production cuts, no field or
optical processes, and no GUI or event output. Five alternating-order repetitions use 200 events per particle
and seeds 12345--12349. Executable and primary hashes are retained. Table~\ref{tab:ft-joints-tracking} reports
median Geant4 beamOn wall times excluding initialization, before/corrected ratios, and paired ratio ranges.
These timings include changed clearances and material, so they are not a measurement of decimation alone.

Median tracking ratios are 0.8568, 0.8541, 0.8465, 0.8920 for electrons, protons, positive pions, and
neutrons, respectively. The corrected median tracking times are 12.1--18.1\% longer than the pre-clearance
baseline. Faster-pair counts are 0/5, 0/5, 0/5, 0/5 in the same order. All 52 final benchmark processes
complete. The 26 corrected processes report no Geant4 exceptions. The baseline overlap diagnostic reports 21
GeomVol1002 blocks. The raw per-seed results retain process CPU, wall time, initialization, and memory
separately from transport.

Loading process wall time changes from 1.279579 to 1.272204 s; CPU changes from 1.190 to 1.230 s. Startup peak
RSS is 180.453 versus 180.188 MiB. The earlier individual-piece benchmarks remain historical measurements of
their stated geometries.

All 9 affected complete run configurations pass construction without Geant4 exceptions. All 17 beamline
configurations pass volume-count construction checks, and the comparator passes 15 native geometry, 15
material, and 14 CAD comparisons. Direct database checks preserve all GEMC2 tables and all GEMC3 physical
rows; every other STL retains its previous hash in all three runtime trees. The historical add-on optimization
remains reproducible using the archived main-cone dependency. These geometry and isolated timing checks do not
establish detector-response equivalence or full CLAS12 performance.

\begin{table}[htbp]
\centering
\caption{Triangle counts and volume changes relative to the previous optimized meshes.}
\label{tab:ft-joints-meshes}
\setlength{\tabcolsep}{5pt}
\begin{tabular}{lrrr}
\toprule
Mesh & Before & Corrected & Volume change (\%) \\
\midrule
Inner extension & 12500 & 11200 & -0.689274 \\
Lead case & 10000 & 16000 & -1.209323 \\
Lead fill & 8000 & 3958 & -0.192478 \\
Tungsten add-on & 12000 & 8356 & -0.374681 \\
Main cone & 5000 & 8546 & -0.118539 \\
\bottomrule
\end{tabular}
\end{table}

\begin{table}[htbp]
\centering
\caption{Placed interface checks; distances sample vertices against the other surface.}
\label{tab:ft-joints-gaps}
\setlength{\tabcolsep}{5pt}
\begin{tabular}{lrr}
\toprule
Interface & Sampled gap (mm) & Contacts \\
\midrule
Shield / inner extension & 0.062045 & 0 \\
Flange / inner extension & 0.053650 & 0 \\
Flange / case & 0.020020 & 0 \\
Case / lead fill & 0.019957 & 0 \\
Lead fill / inner extension & 0.053503 & 0 \\
Main cone / inner extension & 0.054603 & 0 \\
Cone tip / main cone & 0.052979 & 0 \\
Add-on / main cone & 0.094198 & 0 \\
\bottomrule
\end{tabular}
\end{table}

\begin{table}[htbp]
\centering
\caption{Eight-part tracking; five repetitions of 200 events per particle.}
\label{tab:ft-joints-tracking}
\setlength{\tabcolsep}{5pt}
\begin{tabular}{lrrrr}
\toprule
Particle & Before (s) & After (s) & Ratio & Paired range \\
\midrule
Electron & 19.36220 & 22.59820 & 0.8568 & 0.8543--0.8773 \\
Proton & 5.29483 & 6.19904 & 0.8541 & 0.8507--0.8899 \\
Positive pion & 6.94979 & 8.21029 & 0.8465 & 0.8001--0.8923 \\
Neutron & 5.55687 & 6.22962 & 0.8920 & 0.8433--0.9094 \\
\bottomrule
\end{tabular}
\end{table}

\FloatBarrier
\begin{figure}[p]
\centering
\includegraphics[width=\linewidth,height=.80\textheight,keepaspectratio]{figures/ft_joint_sections_1.pdf}
\caption{Matched meridian sections before and after the local clearance corrections. Blue and orange identify
the two
named parts. Coordinates use the common passive cone frame; the shield is translated by its actual offset.
Solid outlines are sections of the stored STLs, with matched axes within each pair.}
\label{fig:ft-joints-1}
\end{figure}
\begin{figure}[p]
\centering
\includegraphics[width=\linewidth,height=.80\textheight,keepaspectratio]{figures/ft_joint_sections_2.pdf}
\caption{Matched meridian sections before and after the local clearance corrections. Blue and orange identify
the two
named parts. Coordinates use the common passive cone frame; the shield is translated by its actual offset.
Solid outlines are sections of the stored STLs, with matched axes within each pair.}
\label{fig:ft-joints-2}
\end{figure}

\FloatBarrier

\section{HTCC CAD cone transfer and validation}
\label{sec:htcc}

The October 19 HTCC update transfers three cured GEMC3 meshes into the GEMC2 clas12Tags repository. These are
the HTCC entry cone, Moller cone, and Moller-cone extension; they are separate from the beamline Moller-cone
studies above. The active source and runtime copies match the GEMC3 assets byte for byte, and the unmodified
GEMC2 engineering inputs are archived separately. Names, Rohacell material, passive 180-degree y rotations,
and all run-dependent placements remain unchanged. The variations default, rga\_spring2018, and rga\_fall2018
use runs 11, 3029, and 4763, respectively; the extension retains its additional 1 mm z shift.

Curing merges vertices, closes mesh boundaries, and repairs orientation. The entry cone is simplified to 6,000
triangles; the Moller cone additionally uses self-intersection removal and finishes at 5,866 triangles. The
already-light extension retains 1,748 triangles. No additional engineering-hole filling or Poisson
reconstruction is applied. Total stored triangles fall from 26,798 to 13,614, a 49.20\% reduction.
Table~\ref{tab:htcc-mesh} lists the topology counts from the same PyMeshLab inspection.


\begin{table}[htbp]
\centering
\caption{HTCC cone topology before and after the transferred curing and simplification.}
\label{tab:htcc-mesh}
\setlength{\tabcolsep}{5pt}
\begin{tabular}{llrrr}
\toprule
Mesh & Input & Triangles & Boundary edges & Non-manifold edges \\
\midrule
\texttt{htccCone} & Raw & 18,818 & 144 & 0 \\
\texttt{htccCone} & Cured & 6,000 & 0 & 0 \\
\texttt{htccMollerCone} & Raw & 6,232 & 0 & 45 \\
\texttt{htccMollerCone} & Cured & 5,866 & 0 & 24 \\
\texttt{htccMollerConeExt} & Raw & 1,748 & 0 & 0 \\
\texttt{htccMollerConeExt} & Cured & 1,748 & 0 & 0 \\
\bottomrule
\end{tabular}
\end{table}

All three cured meshes have zero boundary edges. The entry cone and extension have no non-manifold edges,
while the Moller cone retains 24, reduced from 45 in the raw input. These edge counts depend on the mesh
import and vertex-merging convention and should not be equated to affected-face counts. The largest change
among the six axis-aligned entry-cone bounds is 0.737915 mm. The Moller cone and extension retain their
original bounds exactly. Bounds agreement and zero boundary edges do not establish material-volume or surface
equivalence, and the remaining non-manifold Moller-cone edges are an unresolved topology limitation.

Validation uses GEMC 6.2 with Geant4 10.7.4, a local CCDB SQLite database, no magnetic field, no GUI, no event
output, and zero transported events. Native-plus-CAD initialization uses STD + FTFP\_BERT + Optical. All three
variations complete through both TEXT/CAD and SQLite/CAD, giving six successful initializations. For each
variation, the 537 native TEXT rows match the effective SQLite geometry rows field by field; the three GXML
CAD placements match their effective SQLite CAD rows. Native geometry and materials, CAD placements, and the
SQLite database were retained by the transfer.

The full default SQLite overlap check reports entry-cone contact with htccEntryWindow. An isolated cone/window
fixture uses STD + FTFP\_BERT, CHECK\_OVERLAPS=1000, RANDOM=12345, and the unchanged window geometry and cone
placement. The cured cone produces one GeomVol1002 exception block, with reported penetration of at least
52.4771 micrometres; the matched raw-cone check produces none. The full default check reports at least 52.6119
micrometres for the same pair, approximately 0.053 mm. The two Moller cones report no overlap in that default
assembly check. The entry-window overlap remains unresolved in the transferred mesh; successful construction
does not certify an overlap-free assembly. Sampled penetration reports are not exact overlap depths or
volumes.

The retained zero-event logs verify construction and provide geometry diagnostics. They are not repeated
loading or tracking benchmarks, so no CPU, wall-time, memory, or transport-speed gain is assigned to the
49.20\% triangle reduction. Photon yield, material budget, detector response, and particle navigation
equivalence require separate studies. The HTCC archive retains raw and cured STLs, input fixtures, scripts,
topology hashes, validation logs, and database-row snapshots.

\FloatBarrier

\section{Conclusions}

\begin{enumerate}

\item For LTCC, against the complete original hardware, the optimized CAD bundle contains 93.9\% fewer
triangles and 93.9\% fewer bytes. CAD-only loading CPU time decreases by 92.6\%. This is direct evidence
that the simplified files reduce loading work.

\item Complete LTCC initialization CPU time falls by 63.5\% and startup peak memory falls by 56.8\%. These
gains include both optimized files and the simplified native hierarchy.

\item Nonoptical electron transport is 1.29 times faster; hadron workloads take 1.35--1.53 times as much CPU
time. The separate cone-only measurements quantify the isolated Winston-cone substitution.

\item Full-port optical electron transport requires 1.73 times as much median CPU time in this benchmark. The
increased CPU use may reflect more complete transport after geometry defects are removed: the original model
repeatedly requires corrective navigation pushes, while optimized diagnostic samples contain more tracks and
steps. Lower original CPU time is therefore not sufficient evidence of faster correct transport. Revised
materials and containment also contribute, so the separate causes remain unresolved.

\item The optimized meshes remove 857 exact zero-area triangles and eliminate the 1,714 invalid/rejected-facet
warnings per original import. Navigation reports across the tested workloads fall from 9,546 to 1. Overlap and
containment reports remain, so these gains do not establish overlap-free geometry.

\item Ported TEXT and SQLite geometries agree for all six run-dependent configurations. Warning counts are
reported for the tested construction sample and transport workloads; they do not establish that the complete
CLAS12 geometry is free of overlaps. Detector-response, material-budget, magnetic-field, and full-experiment
comparisons remain outside this study.

\item The separate beamline apex pair contains 45.4\% fewer triangles while preserving topology, sharp hub
features, materials, and placements, with sampled shape differences below 0.05 mm. Its median loading wall
time falls by 15.7\% and startup peak memory by 17.4\%. Transport wall time improves by 1.77--2.01 times in
the first, smaller matched-primary sample. These gains describe the isolated GEMC3 pair and do not establish
full-experiment performance or detector-response equivalence.

\item The downstream tube pair is reduced by 3.75\% in triangles, and Segment 1 is closed by registering and
bridging its missing annular joins while retaining the bore. Import surface-hole warnings are removed, but
overlaps and navigation warnings remain. The measured transport ratios describe this limited workload and do
not establish response equivalence or faster correct transport.

\item The already-coarse adapter and torus exit-shield group permits only a 0.434\% triangle reduction with
protected casing mating faces. The adapter is renamed consistently, and the lead fill is retained exactly.
Its repeated benchmark does not establish substantial performance gains; shape preservation takes priority
over a larger triangle reduction.

\item The FT adjuster trio is replaced by one 48-triangle closed prism, reducing triangles by 98.74\% while
matching the estimated sealed reference union volume within 0.00003\%. The material and common placement are
retained, and internal overlaps disappear in the isolated check. The reference mass remains approximate
because two original shells are open; deliberate material redistribution prevents a claim of
detector-response equivalence.

\item The renamed FT inner shield and flange contains 36.14\% fewer triangles with retained bounds, alloy
material, placement, two closed components, and mounting/bore topology. Its small enclosed pocket is
unchanged, sampled surface differences remain below 0.05 mm, and total signed volume changes by about
-0.000044\%. The isolated benchmark quantifies loading, transport, and diagnostics without establishing
full-experiment response equivalence.

\item Nine renamed FT inner-shield bolts use smooth solid shafts and round heads, reducing triangles by 89.38\%
while retaining each copper volume within 0.0001\%. Shortened tips leave 0.15 mm gaps at the blind-hole
ends; independent collision, containment, and 10,000-point Geant4 checks find no shield interference.
Electron, pion, and neutron transport timings improve in this isolated sample, while proton tracking is
essentially unchanged. These results do not establish full-experiment speedup or response equivalence.

\item The renamed FT outer flange closes twelve bolt tunnels and reduces triangles by 60.79\% while retaining
the
central opening, engineering rim coordinates, bounds, and placement. The authorized hole closure adds
8.296635\% stainless-steel volume. Repeated tracking timings and residual assembly overlaps are recorded
separately; added material prevents attributing timing differences solely to fewer triangles or inferring
response equivalence.

The separate control at fixed material volume gives median tracking speedups of 1.073, 1.016, 1.034, and
1.013 for electrons, protons, positive pions, and neutrons, with the simplified mesh faster in every seed
pair of this larger sample.

\item The renamed FT outer mount repairs four triangular surface gaps and reduces triangles by 11.44\%,
retaining
its existing openings and placement. Protected radial relief removes about 0.24924\% material against the
sealed reference and provides at least 0.03597 mm radial separation from the flange in their common axial
interval. Surface checks and the 10,000-point assembly test find no joint overlap, resolving the earlier
flange/mount interference. Repeated tracking results show modest electron, proton, and pion gains with
slower neutron tracking; unrelated full-beamline diagnostics remain.

\item The final downstream FT inner shield has smooth flange faces with every bolt recess filled, while
retaining
both longitudinal grooves and the central passage. It contains 6,500 triangles, 85.27\% fewer than the raw
input and 51.85\% fewer than the first revision. Recess closure adds 0.436019\% tungsten volume relative to
the prior closed mesh. Upgrade and smooth-reference tracking comparisons are reported separately with clean
isolated transport; the subsequent Moller-cone revision resolves the shield/cone interference.

\item The renamed Moller cone fills its surface openings, retains its beam passage and outer bounds, and
reduces
triangles by 68.32\%. A widened protected counterbore removes about 0.168806\% tungsten relative to the
hole-filled reference and provides at least 0.187891 mm radial shield clearance. The cone/shield assembly
passes the 10,000-point overlap check. Raw-input timing includes solid and navigation warnings; the clean
clearance-reference control separately measures tessellation gains at effectively unchanged material volume.

\item Local relief of both Segment 1 sockets resolves the remaining downstream adapter and tube contacts
without
changing placements or neighboring meshes. Segment 1 has 5,096 triangles, 0.604764\% less steel than its
previous closed mesh, and at least 0.1 mm socket-face clearance. The three-part high-resolution overlap test
and full default beamline construction pass without exceptions. Separate repeated tracking results include
the repair and material change.

\item The beamline runtime database now selects every configuration by run number under \texttt{default}.
Run 16000 loads all 32 geometry volumes, and all 17 configured runs pass construction with unchanged
physical geometry. This selection fix has no independently measured tracking-speed gain.

\item The renamed Moller cone tip has 63.58\% fewer triangles with retained topology, sharp faces, bounds, and
a -0.000486\% material-volume change. Isolated tracking is 1.659, 1.133, 1.177, and 1.159 times faster for
electrons, protons, positive pions, and neutrons in this sample. Its pre-existing main-cone overlap remains in
the four affected complete configurations and is reported separately.

\item The renamed FT flange retains its opening, flat faces, bounds, and topology with 540 triangles, 2.17\%
fewer than the already-coarse original. Steel volume changes by -0.000109\%, and sampled surface differences
remain below 0.05 mm. Median tracking ratios remain within 0.18\% of one, providing no meaningful speed-gain
evidence. Existing neighboring assembly defects remain and are reported separately.

\item The renamed inner Moller-cone extension restores its missing radial-port wall and reduces triangles by
30.77\% while retaining its passages, bounds, and protected faces. Decimation changes the repaired reference
volume by +0.000009303\%, with sampled deviations below 0.016 mm. Repaired-reference tracking ratios are
1.261, 1.086, 1.144, 1.095 for electrons, protons, positive pions, and neutrons. The original is open and
produces navigation warnings; those raw timings are reported separately. Existing assembly overlaps remain.

\item The moller-cone extension lead fill retains its mating surface and openings with 24.49\% fewer
triangles. Median isolated tracking ratios are 1.0826, 1.0795, 1.1226, 1.0882 for electrons, protons, positive
pions, and neutrons. Every paired measurement improves in this sample. Existing assembly overlaps remain.

\item The moller-cone extension lead case retains its mating surface and openings with 42.40\% fewer
triangles. Median isolated tracking ratios are 1.0529, 1.0377, 1.0098, 1.0198 for electrons, protons, positive
pions, and neutrons. Some paired ranges include one; those particle samples do not establish a consistent
gain. Existing assembly overlaps remain.

\item The renamed Moller-cone extension tungsten add-on restores its missing radial-port wall and reduces
triangles by 56.80\% while retaining its passages, bounds, and protected faces. Decimation changes the
repaired reference volume by -0.0000046193\%, with sampled deviations below 0.013 mm. Repaired-reference
tracking ratios are 1.647, 1.548, 1.642, 1.535 for electrons, protons, positive pions, and neutrons. The
original is open and produces navigation warnings; those raw timings are reported separately. Existing
assembly overlaps remain.

\item Local FT/cone joint reliefs clear all eight checked interfaces while retaining the shield, flange, and
cone tip. Five meshes lose 0.119--1.209\% material by volume. All affected complete configurations construct
without exceptions, and final corrected tracking runs are clean. Isolated tracking ratios 0.8568, 0.8541,
0.8465, 0.8920 describe the eight-part clearance revision; they are not pure decimation gains.

\item The three HTCC cones have 49.20\% fewer triangles, and the entry cone's 144 boundary edges are closed.
All six TEXT/CAD and SQLite/CAD initialization checks and the three native/CAD database comparisons pass. The
Moller cone retains 24 non-manifold edges, and the cured entry cone introduces an observed contact of
approximately 0.053 mm with the entry window. These unresolved defects preclude a claim of complete geometry
validation; no performance or detector-response gain has been measured for this transfer.

\end{enumerate}

\FloatBarrier

\section{Reproducibility}

The companion data archive retains metadata and per-run JSON results for each workload, all raw logs and time
reports, the geometry validation summary, and matched primary-input hashes. The repository benchmark script
runs offline with the local CCDB database. The retained metadata provides the exact executable hash; use that
binary and the original revision for exact reproduction. Timing results are specific to this machine and
workload. The retained dataset is cited in Ref.~\cite{dataset}.

The separate beamline apex dataset retains both original and optimized two-mesh fixtures, all five matched
primary files per particle, every log and time report, mesh-quality statistics, the builder, benchmark
script, and executable hashes. It is stored in the companion beamline archive; the executable is not
redistributed. The beamline benchmark needs no CCDB connection.

The separate downstream and adapter/exit-shield archives retain the corresponding fixtures, matched
primaries, raw logs, time reports, mesh checks, executable hashes, and reproduction scripts. Each beamline
workload has its own archive and data directory; the unchanged lead-fill geometry is included in both exit-
group fixtures. These archives do not redistribute the executables.

The FT adjuster archive additionally retains the original three-piece fixture, the one-solid replacement,
material reference checks, and the model builder. The reference mass uses sealed imperfect shells and
overlap subtraction, so it must not be treated as a certified engineering mass.

The inner-shield archive retains the original engineering mesh, the optimized single-mesh fixture, the
unchanged enclosed-pocket geometry, component volume checks, matched primaries, logs, and timing reports.
The original stem is retained in the engineering archive while runtime placements and both fixtures use the
new name.

The bolt archive retains the nine original engineering meshes, volume-matched replacements, unchanged shield
used for clearance, surface collision and containment checks, high-resolution Geant4 overlap logs, and all
matched transport inputs and timings. Earlier tracking results remain in the five preceding beamline
sections.

The outer-flange archive retains the original engineering mesh, the hole-filled reference, exact-rim
retriangulation, matched primaries, transport and loading measurements, and original/optimized assembly
fixtures. Only the fixed-executable benchmark is reported; discarded runs with a changing executable are
excluded. The snapshot hash and hole-removal utility source identify the measured implementation without
distributing a binary.

The outer-flange archive also includes a separate material-matched control, comparing the hole-filled
reference with the simplified mesh for 2,000 events per run. The new benchmark option selects the filled
baseline explicitly; geometry and material records, matched primary hashes, paired seed ratios, and both
complete fixtures are retained. The earlier drilled-versus-filled results remain independently documented.

The outer-mount archive retains the open engineering mesh, sealed material reference, protected joint
relief, unchanged mating flange, and full original/optimized assembly fixtures. Clipped triangle bounds and
VTK checks verify placed clearance independently of Geant4 sampling. Tracking times, matched primaries,
executable hashes, all raw diagnostic logs, and byte-for-byte reproduction evidence are included.

The FT inner-shield archive retains the damaged engineering input, hole-filled reconstructed reference,
final protected surface, independent feature probes, and assembly checks with the unchanged cone.
Engineering and reconstructed-reference transport comparisons have separate matched-primary records, timing
reports, and diagnostic logs. Original mass is undefined; nominal profile agreement and subsequent
decimation checks are distinguished.

A supplemental FT inner-shield archive retains the smooth-flange revision, its previous retained-recess
mesh, and two independent repeated tracking comparisons. Complete flat annulus areas and dense near-face
containment probes verify recess closure. The original study archive remains unchanged, preserving its
historical code and measurements.

The Moller-cone archive retains the original open mesh, both closed references, the final cone and unchanged
FT inner shield, all nine GXML inputs, and full benchmark and validation evidence. Original-input timings
include source solid and navigation warnings. A separate clean-reference control measures decimation after
the authorized hole fill and clearance relief.

A second supplemental archive retains the downstream socket revision, the previous three-part assembly, all
repeated tracking runs, clipped-surface support checks, and full construction logs. Both neighboring meshes
and all CAD placement rows remain unchanged. Previous study archives retain their original bytes and
historical overlap measurements.


The separate cone-tip archive retains the clean engineering input, optimized tip, unchanged main cone,
both isolated and assembly fixtures, all repeated timing logs, matched primaries, and the four affected
complete configurations. Mesh and database checks distinguish the rename and decimation from pre-existing
assembly defects. The run-5874 benchmark uses the default variation; its executable hash and source
snapshots are retained without redistributing the executable. All earlier archives remain unchanged.


The separate FT-flange archive retains the original and optimized flange, both neighboring meshes,
all repeated timing logs and primary files, all four complete configurations, and direct mesh/database
checks. Run 5874 uses the default variation in both isolated fixtures. The small triangle reduction
shows no meaningful tracking gain, and existing assembly defects remain. Earlier archives are unchanged.


The inner-extension archive retains the open source, closed reference, optimized mesh, matched primary
files, and both complete repeated tracking comparisons. A 75-triangle port-wall repair preserves all
original vertices. Feature probes and protected-face checks verify retained passages and engineering
features. Placed assembly diagnostics, all four affected complete configurations, database backups,
source snapshots, and executable hashes are retained. Existing assembly overlaps remain; raw defective
input timings and clean repaired-reference results are separate. Earlier archives remain unchanged.


The separate lead-fill and lead-case extension archives retain matched five-seed timing inputs, raw process
reports, both database fixtures, original and optimized meshes, assembly checks, all four complete affected
configurations, and code/provenance snapshots. Both studies select run 5874 with the default variation.
Existing assembly overlaps remain. Earlier archives are unchanged.


The tungsten add-on archive retains the open source, closed reference, optimized mesh, matched primary
files, and both complete repeated tracking comparisons. A 68-triangle inclined-port wall repair preserves all
original vertices. Feature probes and protected-face checks verify retained passages and engineering
features. Placed assembly diagnostics, all four affected complete configurations, database backups,
source snapshots, and executable hashes are retained. Existing assembly overlaps remain; raw defective
input timings and clean repaired-reference results are separate. Earlier archives remain unchanged.


The FT/cone clearance archive retains the previous and corrected eight-part assembly, all final matched timing
inputs and process reports, the excluded cap trial, independent interface checks, all affected complete run
configurations, and code/provenance snapshots. Archived optimized inputs make the final pass independent of
later changes to individual builders. Earlier archives are unchanged.

The separate HTCC archive retains the October 19 raw and cured cones, the six zero-event initialization
fixtures, three run-dependent native/CAD row snapshots, the full default overlap log, and the matched
raw/cured entry-window fixtures. Topology and STL hashes identify the exact transferred meshes. These checks
include no timing benchmark or response-equivalence study. The entry-window overlap and non-manifold
Moller-cone edges remain unresolved. Earlier archives are unchanged.


\FloatBarrier
\begin{thebibliography}{9}
\bibitem{baseline}
GEMC clas12Tags, main revision \texttt{ea3625f}: original LTCC geometry and CAD data;
\texttt{cadTempRemoved} hardware restored using the legacy \texttt{ltcc\_frame.pl} placements.
\bibitem{optimized}
GEMC clas12Tags, meson revision \texttt{4c7c6af}, ``new optimized STL for ltcc and new mother
volume for each ltcc sector'': optimized meshes, CAD placements, mother geometry, and effective nose material.
\bibitem{gemc}
M. Ungaro et al., ``The CLAS12 Geant4 simulation'', Nuclear Instruments and Methods A 959 (2020) 163422.
\bibitem{cad}
R. De Vita and M. Ungaro, ``CLAS12-note 2017-017: Importing CLAS12 CAD Models of Target and Beamline''.
\bibitem{dataset}
Companion datasets, LTCC and beamline CAD optimization studies, October 1, October 12--16, and October 19,
2026:
measured configurations, warnings, timing, memory, hashes, and validation results;
\texttt{study\_archive.zip}, the sixteen beamline study archives, and the supplemental smooth-flange and
downstream-joint
revision archives and the HTCC cone validation archive accompanying this note.
\end{thebibliography}

\end{document}
